\documentclass[10pt]{article}
\usepackage[utf8]{inputenc}
\usepackage[spanish,es-tabla]{babel}
\usepackage[margin=0.8cm,top=0.6cm,bottom=0.6cm]{geometry}
\usepackage{helvet}
\renewcommand{\familydefault}{\sfdefault}
\usepackage{microtype}
\usepackage{graphicx}
\usepackage{xcolor}
\usepackage{tikz}
\usepackage{enumitem}
\usepackage{fontawesome5}
\usepackage{tcolorbox}
\tcbuselibrary{skins}

\pagestyle{empty}

% Indicamos que las plantillas se cargan como módulos
\def\isSubFile{1}
% ==============================================================================
% PLANTILLA_ALUMNO.TEX — MACRO Y DEFINICIÓN DE FICHA PARA ALUMNADO (NEAE)
% Asignatura: Metodologías de Desarrollo Ágil (MDA) — Universidad de Granada
% Proyecto AUTOVIDA (26-175) | Reto CeRA — Agenda Accesible PTVAL
% ==============================================================================
% Campos canónicos: 9 (Sin Occupation, Education ni Company Information)
% Uso: % ==============================================================================
% PLANTILLA_ALUMNO.TEX — MACRO Y DEFINICIÓN DE FICHA PARA ALUMNADO (NEAE)
% Asignatura: Metodologías de Desarrollo Ágil (MDA) — Universidad de Granada
% Proyecto AUTOVIDA (26-175) | Reto CeRA — Agenda Accesible PTVAL
% ==============================================================================
% Campos canónicos: 9 (Sin Occupation, Education ni Company Information)
% Uso: % ==============================================================================
% PLANTILLA_ALUMNO.TEX — MACRO Y DEFINICIÓN DE FICHA PARA ALUMNADO (NEAE)
% Asignatura: Metodologías de Desarrollo Ágil (MDA) — Universidad de Granada
% Proyecto AUTOVIDA (26-175) | Reto CeRA — Agenda Accesible PTVAL
% ==============================================================================
% Campos canónicos: 9 (Sin Occupation, Education ni Company Information)
% Uso: \input{plantilla_alumno.tex} y luego invocar \fichaAlumno[...] o \fichaAlumnoBlanco
% ==============================================================================

\ifdefined\isSubFile
    % Si se incluye desde otro documento, los paquetes ya están cargados
\else
    \documentclass[10pt]{article}
    \usepackage[utf8]{inputenc}
    \usepackage[spanish,es-tabla]{babel}
    \usepackage[margin=0.8cm,top=0.6cm,bottom=0.6cm]{geometry}
    \usepackage{helvet}
    \renewcommand{\familydefault}{\sfdefault}
    \usepackage{microtype}
    \usepackage{graphicx}
    \usepackage{xcolor}
    \usepackage{tikz}
    \usepackage{enumitem}
    \usepackage{fontawesome5}
    \usepackage{tcolorbox}
    \tcbuselibrary{skins}
    \pagestyle{empty}
\fi

% Colores del sistema
\providecolor{primaryColor}{RGB}{20, 65, 105}     % Azul corporativo
\providecolor{secondaryColor}{RGB}{35, 115, 170}  % Azul medio accesible
\providecolor{tealColor}{RGB}{20, 130, 140}       % Verde azulado
\providecolor{greenColor}{RGB}{30, 135, 75}       % Verde bosque
\providecolor{accentColor}{RGB}{215, 105, 25}     % Naranja cálido
\providecolor{purpleColor}{RGB}{120, 60, 150}     % Púrpura suave
\providecolor{dangerColor}{RGB}{185, 45, 35}      % Rojo criticidad
\providecolor{bgCard}{RGB}{252, 253, 255}         % Fondo limpio
\providecolor{borderCard}{RGB}{215, 225, 235}     % Borde sutil
\providecolor{textDark}{RGB}{30, 40, 50}          % Texto principal oscuro
\providecolor{textMuted}{RGB}{95, 110, 125}       % Texto secundario

% Configuración de listas
\setlist[itemize]{leftmargin=3.0mm, itemsep=0.4pt, parsep=0pt, topsep=0.4pt, partopsep=0pt}

% Estilos de caja fija para Alumno
\tcbset{
    cardAlumnoBox/.style={
        enhanced,
        boxrule=0.65pt,
        arc=1.8mm,
        height=32.0mm,
        left=5.0pt, right=5.0pt, top=1.8pt, bottom=1.8pt,
        boxsep=0pt,
        colback=bgCard,
        colframe=borderCard,
        fonttitle=\bfseries\small,
        fontupper=\fontsize{8.5pt}{10.0pt}\selectfont,
    },
    cardAlumnoWideBox/.style={
        enhanced,
        boxrule=0.65pt,
        arc=1.8mm,
        height=22.5mm,
        left=6.0pt, right=6.0pt, top=2.0pt, bottom=2.0pt,
        boxsep=0pt,
        colback=bgCard,
        colframe=borderCard,
        fonttitle=\bfseries\small,
    }
}

% Macro para badges visuales (si no existe)
\providecommand{\badge}[2]{%
    \tikz[baseline=(b.base)] \node[fill=#1!12, text=#1!90!black, draw=#1!35, line width=0.4pt, rounded corners=1.2mm, inner sep=2.0pt, font=\bfseries\scriptsize] (b) {#2};%
}

% Definición de claves con pgfkeys para \fichaAlumno
\pgfkeys{
    /fichaAlumno/.is family, /fichaAlumno,
    default/.style = {
        bannerTitulo = {\faUserGraduate\ FICHA DE PERSONA: ALUMNADO (METODOLOGÍA DCU)},
        subtitulo = {Proyecto AUTOVIDA (26-175) \textbar\ Reto CeRA — Agenda Accesible PTVAL \textbar\ Asignatura MDA (UGR)},
        equipo = {\textbf{Equipo CeRA-MDA} (5 integrantes)\\\ Course 2026/2027 \textbar\ \textbf{Práctica 1}},
        nombre = {[Nombre y Apellidos del Alumno]},
        nombreCorto = {[ALUMNO]},
        edad = {[Edad]},
        icono = {\faUserGraduate},
        foto = {},
        badgeRol = {[Perfil NEAE / Aula]},
        badgeTarea = {Tareas Rotativas},
        centro = {CEE Purísima Concepción},
        cita = {[Cita textual representativa que exprese cómo vive su rutina, qué le ayuda a anticipar y qué le genera ansiedad o rechazo].},
        demographic = {
            \item \textbf{Edad y Género (Age / Gender):} [[Edad] años, [Género]].
            \item \textbf{Ubicación (Location):} [[Granada capital / Municipio]].
            \item \textbf{Perfil diagnóstico:} [[Diagnóstico cognitivo, motor o sensorial]].
        },
        computerExp = {
            \item \textbf{Habilidades informáticas:} [[Nivel táctil: autónomo, asistido]].
            \item \textbf{Años de experiencia:} [[Años de contacto con tablets]].
            \item \textbf{Hábitos de uso:} [[Apps frecuentes: vídeos, pictogramas, juegos]].
            \item \textbf{Autonomía digital:} [[Baja / Media / Asistida]].
        },
        productExp = {
            \item \textbf{Productos similares:} [[PECS en papel, velcro, apps previas]].
            \item \textbf{Limitaciones actuales:} [[Pérdida de material, rigidez, frustración]].
            \item \textbf{Tendencias de uso:} [[Receptividad a imágenes reales vs pictogramas]].
        },
        tasks = {
            \item \textbf{Tareas Primarias:} 1) [[Login accesible]] $\rightarrow$ 2) [[Ver tarea actual]] $\rightarrow$ 3) [[Ejecutar pasos guiados]] $\rightarrow$ 4) [[Confirmar finalización]].
            \item \textbf{Tareas Secundarias:} [[Ver siguiente tarea / pedir ayuda]].
        },
        domainKnowledge = {
            \item \textbf{Comprensión del área:} [[Orientación en el colegio, aulas y profesores]].
            \item \textbf{Comprensión temporal:} [[Visual (disco/reloj de arena), no numérica]].
            \item \textbf{Destreza procedimental:} [[Grado de autonomía en la tarea asignada]].
        },
        technology = {
            \item \textbf{Hardware:} [[Tablet Android / pulsador externo]].
            \item \textbf{Software:} [[Android / síntesis de voz]].
            \item \textbf{Herramientas habituales:} [[Conexión Wi-Fi / atril]].
        },
        attitudes = {
            \item \textbf{Preferencias:} [[Refuerzos sensoriales favoritos, limpieza visual]].
            \item \textbf{Miedo a la tecnología:} [[Reacción a errores, sonidos o bloqueos]].
            \item \textbf{Valores:} [[Sentimiento de logro, orgullo por autonomía]].
        },
        learningStyle = {
            \item \textbf{Perfil de aprendizaje:} [[\textbf{Visual} / \textbf{Auditivo} / \textbf{Kinestésico}]].
            \item \textbf{Canales óptimos:} [[Fotos reales / Pictogramas ARASAAC / Vídeos]].
            \item \textbf{Metodología:} [[Modelado guiado, repetición, ensayo y éxito]].
        },
        criticality = {
            \item \textbf{Consecuencias de un error:} \textit{Emocional:} [[¿Bloqueo, crisis de frustración, rechazo?]] \textit{Operativo:} [[¿Desorientación, interrupción de la tarea?]]
            \item \textbf{Tolerancia del software:} [[Acciones no destructivas, confirmación siempre asistida, botón SOS visible]].
        },
        escenario = {[[Narrar en 3-4 líneas una situación real: a qué hora usa la agenda, cómo anticipa su rutina con apoyos visuales y cómo supera el momento crítico logrando autonomía sin estrés.]]},
        badgesAjustes = {
            \badge{primaryColor}{\faKey\ Login: [[Método]]} \hspace{2mm}
            \badge{secondaryColor}{\faIcons\ [[ARASAAC / Fotos]]} \hspace{2mm}
            \badge{accentColor}{\faClock\ [[Temporizador]]} \hspace{2mm}
            \badge{purpleColor}{\faVolumeUp\ [[Voz TTS]]}%
        }
    },
    bannerTitulo/.store in = \fAlBannerTitulo,
    subtitulo/.store in = \fAlSubtitulo,
    equipo/.store in = \fAlEquipo,
    nombre/.store in = \fAlNombre,
    nombreCorto/.store in = \fAlNombreCorto,
    edad/.store in = \fAlEdad,
    icono/.store in = \fAlIcono,
    foto/.store in = \fAlFoto,
    badgeRol/.store in = \fAlBadgeRol,
    badgeTarea/.store in = \fAlBadgeTarea,
    centro/.store in = \fAlCentro,
    cita/.store in = \fAlCita,
    demographic/.store in = \fAlDemographic,
    computerExp/.store in = \fAlComputerExp,
    productExp/.store in = \fAlProductExp,
    tasks/.store in = \fAlTasks,
    domainKnowledge/.store in = \fAlDomainKnowledge,
    technology/.store in = \fAlTechnology,
    attitudes/.store in = \fAlAttitudes,
    learningStyle/.store in = \fAlLearningStyle,
    criticality/.store in = \fAlCriticality,
    escenario/.store in = \fAlEscenario,
    badgesAjustes/.store in = \fAlBadgesAjustes,
}

% COMANDO PRINCIPAL \fichaAlumno
\newcommand{\fichaAlumno}[1][]{%
    \begingroup
    \pgfkeys{/fichaAlumno, default, #1}%
    % 1. BANNER SUPERIOR
    \begin{tcolorbox}[
        enhanced,
        colback=primaryColor,
        colframe=primaryColor,
        arc=1.6mm,
        left=7pt, right=7pt, top=2.5pt, bottom=2.5pt,
        boxsep=0pt, coltext=white
    ]
        \begin{minipage}[c]{0.70\textwidth}
            {\normalsize\textbf{\fAlBannerTitulo}}\\[0.2mm]
            {\tiny \fAlSubtitulo}
        \end{minipage}\hfill
        \begin{minipage}[c]{0.29\textwidth}
            \raggedleft\tiny
            \fAlEquipo
        \end{minipage}
    \end{tcolorbox}
    \vspace{0.8mm}
    % 2. TARJETA DE IDENTIDAD DEL ALUMNO
    \begin{tcolorbox}[
        enhanced,
        colback=white,
        colframe=primaryColor!65,
        arc=2.0mm,
        boxrule=0.8pt,
        left=6pt, right=6pt, top=3.0pt, bottom=3.0pt,
        boxsep=0pt,
        drop shadow=black!5
    ]
        \begin{minipage}[c]{0.10\textwidth}
            \centering
            \begin{tikzpicture}
                \ifx\fAlFoto\empty
                    \draw[fill=primaryColor!8, draw=primaryColor!30, line width=0.5pt, rounded corners=2mm] (0,0) rectangle (1.4, 1.55);
                    \node at (0.7, 0.95) {\color{primaryColor}\scalebox{1.5}{\fAlIcono}};
                    \node[align=center] at (0.7, 0.28) {\tiny\textbf{\color{primaryColor}\fAlNombreCorto}\\\fontsize{4pt}{5pt}\selectfont\color{textMuted}\fAlEdad};
                \else
                    \begin{scope}
                        \clip[rounded corners=2mm] (0,0) rectangle (1.45, 1.55);
                        \node[inner sep=0pt] at (0.725, 0.775) {\includegraphics[width=1.45cm,height=1.55cm,keepaspectratio=false]{\fAlFoto}};
                    \end{scope}
                    \draw[draw=primaryColor!50, line width=0.8pt, rounded corners=2mm] (0,0) rectangle (1.45, 1.55);
                \fi
            \end{tikzpicture}
        \end{minipage}\hfill
        \begin{minipage}[c]{0.88\textwidth}
            \begin{minipage}[b]{0.68\textwidth}\raggedright
                {\large\textbf{\color{primaryColor}\fAlNombre}} \hspace{2mm}
                \badge{secondaryColor}{\fAlBadgeRol} \badge{accentColor}{\fAlBadgeTarea}
            \end{minipage}\hfill
            \begin{minipage}[b]{0.31\textwidth}
                \raggedleft\small\textbf{\color{primaryColor}\faSchool\ \fAlCentro}
            \end{minipage}\\[0.6mm]
            \tikz[baseline] \node[fill=primaryColor!5, draw=primaryColor!25, rounded corners=1.2mm, inner sep=2.5pt, text width=0.98\linewidth] {
                \small\itshape\color{primaryColor}
                \faQuoteLeft\ ``\fAlCita''\ \faQuoteRight
            };
        \end{minipage}
    \end{tcolorbox}
    \vspace{0.8mm}
    % 3. LAS 9 CAJAS DE ALUMNOS (4 FILAS X 2 COLS + 1 FILA COMPLETA)
    \noindent\begin{minipage}[t]{0.485\textwidth}
        \begin{tcolorbox}[cardAlumnoBox, colframe=primaryColor!50, coltitle=primaryColor, title={\faIdCard\ 1. Demographic characteristics}]
            \fontsize{8.2pt}{9.8pt}\selectfont
            \begin{itemize}
                \fAlDemographic
            \end{itemize}
        \end{tcolorbox}
    \end{minipage}\hfill
    \begin{minipage}[t]{0.485\textwidth}
        \begin{tcolorbox}[cardAlumnoBox, colframe=secondaryColor!60, coltitle=secondaryColor!90!black, title={\faLaptop\ 2. Computer experience}]
            \fontsize{8.2pt}{9.8pt}\selectfont
            \begin{itemize}
                \fAlComputerExp
            \end{itemize}
        \end{tcolorbox}
    \end{minipage}
    \vspace{0.8mm}
    \noindent\begin{minipage}[t]{0.485\textwidth}
        \begin{tcolorbox}[cardAlumnoBox, colframe=secondaryColor!60, coltitle=secondaryColor!90!black, title={\faSearchPlus\ 3. Specific product experience}]
            \fontsize{8.2pt}{9.8pt}\selectfont
            \begin{itemize}
                \fAlProductExp
            \end{itemize}
        \end{tcolorbox}
    \end{minipage}\hfill
    \begin{minipage}[t]{0.485\textwidth}
        \begin{tcolorbox}[cardAlumnoBox, colframe=greenColor!65, coltitle=greenColor!90!black, title={\faTasks\ 4. Tasks (Primary and secondary)}]
            \fontsize{8.2pt}{9.8pt}\selectfont
            \begin{itemize}
                \fAlTasks
            \end{itemize}
        \end{tcolorbox}
    \end{minipage}
    \vspace{0.8mm}
    \noindent\begin{minipage}[t]{0.485\textwidth}
        \begin{tcolorbox}[cardAlumnoBox, colframe=accentColor!65, coltitle=accentColor!90!black, title={\faBrain\ 5. Domain knowledge}]
            \fontsize{8.2pt}{9.8pt}\selectfont
            \begin{itemize}
                \fAlDomainKnowledge
            \end{itemize}
        \end{tcolorbox}
    \end{minipage}\hfill
    \begin{minipage}[t]{0.485\textwidth}
        \begin{tcolorbox}[cardAlumnoBox, colframe=tealColor!60, coltitle=tealColor!90!black, title={\faTools\ 6. Technology available}]
            \fontsize{8.2pt}{9.8pt}\selectfont
            \begin{itemize}
                \fAlTechnology
            \end{itemize}
        \end{tcolorbox}
    \end{minipage}
    \vspace{0.8mm}
    \noindent\begin{minipage}[t]{0.485\textwidth}
        \begin{tcolorbox}[cardAlumnoBox, colframe=purpleColor!60, coltitle=purpleColor!90!black, title={\faHeart\ 7. Attitudes and values}]
            \fontsize{8.2pt}{9.8pt}\selectfont
            \begin{itemize}
                \fAlAttitudes
            \end{itemize}
        \end{tcolorbox}
    \end{minipage}\hfill
    \begin{minipage}[t]{0.485\textwidth}
        \begin{tcolorbox}[cardAlumnoBox, colframe=purpleColor!60, coltitle=purpleColor!90!black, title={\faEye\ 8. Learning style}]
            \fontsize{8.2pt}{9.8pt}\selectfont
            \begin{itemize}
                \fAlLearningStyle
            \end{itemize}
        \end{tcolorbox}
    \end{minipage}
    \vspace{0.8mm}
    % FILA 5: Punto 9 (ANCHO COMPLETO — CRITICALITY OF ERRORS)
    \noindent\begin{minipage}[t]{\textwidth}
        \begin{tcolorbox}[cardAlumnoWideBox, colframe=dangerColor!65, coltitle=dangerColor!90!black, title={\faExclamationTriangle\ 9. Criticality of errors (Consecuencias de un error y tolerancia del software)}]
            \fontsize{8.2pt}{9.8pt}\selectfont
            \begin{itemize}
                \fAlCriticality
            \end{itemize}
        \end{tcolorbox}
    \end{minipage}
    \vspace{0.8mm}
    % 4. BARRA INFERIOR DE INTEGRACIÓN CeRA
    \begin{tcolorbox}[
        enhanced,
        colback=primaryColor!4,
        colframe=primaryColor!45,
        arc=1.6mm,
        left=6pt, right=6pt, top=2.5pt, bottom=2.5pt,
        boxsep=0pt
    ]
        \scriptsize
        \textbf{\color{primaryColor}\faMapSigns\ Escenario de Uso Contextual (Día típico):} \fAlEscenario\\
        \vspace{0.3mm}
        \hrule height 0.35pt \relax
        \vspace{0.5mm}
        \centerline{%
            \textbf{\color{primaryColor}\faSlidersH\ Ajustes UI Accesible (CeRA):} \hspace{2mm}
            \fAlBadgesAjustes%
        }
    \end{tcolorbox}
    \newpage
    \endgroup
}

% COMANDO PARA PLANTILLA EN BLANCO DE ALUMNADO
\newcommand{\fichaAlumnoBlanco}{%
    \fichaAlumno[%
        bannerTitulo={\faUserEdit\ FICHA DE PERSONA: ALUMNADO \textbar\ PLANTILLA DE TRABAJO},
        cita={[Cita textual que exprese cómo vive su rutina, qué le ayuda a anticipar y qué le genera ansiedad o rechazo].}%
    ]%
}

\ifdefined\isSubFile
    \expandafter\endinput
\fi

\begin{document}
\fichaAlumnoBlanco
\end{document}
 y luego invocar \fichaAlumno[...] o \fichaAlumnoBlanco
% ==============================================================================

\ifdefined\isSubFile
    % Si se incluye desde otro documento, los paquetes ya están cargados
\else
    \documentclass[10pt]{article}
    \usepackage[utf8]{inputenc}
    \usepackage[spanish,es-tabla]{babel}
    \usepackage[margin=0.8cm,top=0.6cm,bottom=0.6cm]{geometry}
    \usepackage{helvet}
    \renewcommand{\familydefault}{\sfdefault}
    \usepackage{microtype}
    \usepackage{graphicx}
    \usepackage{xcolor}
    \usepackage{tikz}
    \usepackage{enumitem}
    \usepackage{fontawesome5}
    \usepackage{tcolorbox}
    \tcbuselibrary{skins}
    \pagestyle{empty}
\fi

% Colores del sistema
\providecolor{primaryColor}{RGB}{20, 65, 105}     % Azul corporativo
\providecolor{secondaryColor}{RGB}{35, 115, 170}  % Azul medio accesible
\providecolor{tealColor}{RGB}{20, 130, 140}       % Verde azulado
\providecolor{greenColor}{RGB}{30, 135, 75}       % Verde bosque
\providecolor{accentColor}{RGB}{215, 105, 25}     % Naranja cálido
\providecolor{purpleColor}{RGB}{120, 60, 150}     % Púrpura suave
\providecolor{dangerColor}{RGB}{185, 45, 35}      % Rojo criticidad
\providecolor{bgCard}{RGB}{252, 253, 255}         % Fondo limpio
\providecolor{borderCard}{RGB}{215, 225, 235}     % Borde sutil
\providecolor{textDark}{RGB}{30, 40, 50}          % Texto principal oscuro
\providecolor{textMuted}{RGB}{95, 110, 125}       % Texto secundario

% Configuración de listas
\setlist[itemize]{leftmargin=3.0mm, itemsep=0.4pt, parsep=0pt, topsep=0.4pt, partopsep=0pt}

% Estilos de caja fija para Alumno
\tcbset{
    cardAlumnoBox/.style={
        enhanced,
        boxrule=0.65pt,
        arc=1.8mm,
        height=32.0mm,
        left=5.0pt, right=5.0pt, top=1.8pt, bottom=1.8pt,
        boxsep=0pt,
        colback=bgCard,
        colframe=borderCard,
        fonttitle=\bfseries\small,
        fontupper=\fontsize{8.5pt}{10.0pt}\selectfont,
    },
    cardAlumnoWideBox/.style={
        enhanced,
        boxrule=0.65pt,
        arc=1.8mm,
        height=22.5mm,
        left=6.0pt, right=6.0pt, top=2.0pt, bottom=2.0pt,
        boxsep=0pt,
        colback=bgCard,
        colframe=borderCard,
        fonttitle=\bfseries\small,
    }
}

% Macro para badges visuales (si no existe)
\providecommand{\badge}[2]{%
    \tikz[baseline=(b.base)] \node[fill=#1!12, text=#1!90!black, draw=#1!35, line width=0.4pt, rounded corners=1.2mm, inner sep=2.0pt, font=\bfseries\scriptsize] (b) {#2};%
}

% Definición de claves con pgfkeys para \fichaAlumno
\pgfkeys{
    /fichaAlumno/.is family, /fichaAlumno,
    default/.style = {
        bannerTitulo = {\faUserGraduate\ FICHA DE PERSONA: ALUMNADO (METODOLOGÍA DCU)},
        subtitulo = {Proyecto AUTOVIDA (26-175) \textbar\ Reto CeRA — Agenda Accesible PTVAL \textbar\ Asignatura MDA (UGR)},
        equipo = {\textbf{Equipo CeRA-MDA} (5 integrantes)\\\ Course 2026/2027 \textbar\ \textbf{Práctica 1}},
        nombre = {[Nombre y Apellidos del Alumno]},
        nombreCorto = {[ALUMNO]},
        edad = {[Edad]},
        icono = {\faUserGraduate},
        foto = {},
        badgeRol = {[Perfil NEAE / Aula]},
        badgeTarea = {Tareas Rotativas},
        centro = {CEE Purísima Concepción},
        cita = {[Cita textual representativa que exprese cómo vive su rutina, qué le ayuda a anticipar y qué le genera ansiedad o rechazo].},
        demographic = {
            \item \textbf{Edad y Género (Age / Gender):} [[Edad] años, [Género]].
            \item \textbf{Ubicación (Location):} [[Granada capital / Municipio]].
            \item \textbf{Perfil diagnóstico:} [[Diagnóstico cognitivo, motor o sensorial]].
        },
        computerExp = {
            \item \textbf{Habilidades informáticas:} [[Nivel táctil: autónomo, asistido]].
            \item \textbf{Años de experiencia:} [[Años de contacto con tablets]].
            \item \textbf{Hábitos de uso:} [[Apps frecuentes: vídeos, pictogramas, juegos]].
            \item \textbf{Autonomía digital:} [[Baja / Media / Asistida]].
        },
        productExp = {
            \item \textbf{Productos similares:} [[PECS en papel, velcro, apps previas]].
            \item \textbf{Limitaciones actuales:} [[Pérdida de material, rigidez, frustración]].
            \item \textbf{Tendencias de uso:} [[Receptividad a imágenes reales vs pictogramas]].
        },
        tasks = {
            \item \textbf{Tareas Primarias:} 1) [[Login accesible]] $\rightarrow$ 2) [[Ver tarea actual]] $\rightarrow$ 3) [[Ejecutar pasos guiados]] $\rightarrow$ 4) [[Confirmar finalización]].
            \item \textbf{Tareas Secundarias:} [[Ver siguiente tarea / pedir ayuda]].
        },
        domainKnowledge = {
            \item \textbf{Comprensión del área:} [[Orientación en el colegio, aulas y profesores]].
            \item \textbf{Comprensión temporal:} [[Visual (disco/reloj de arena), no numérica]].
            \item \textbf{Destreza procedimental:} [[Grado de autonomía en la tarea asignada]].
        },
        technology = {
            \item \textbf{Hardware:} [[Tablet Android / pulsador externo]].
            \item \textbf{Software:} [[Android / síntesis de voz]].
            \item \textbf{Herramientas habituales:} [[Conexión Wi-Fi / atril]].
        },
        attitudes = {
            \item \textbf{Preferencias:} [[Refuerzos sensoriales favoritos, limpieza visual]].
            \item \textbf{Miedo a la tecnología:} [[Reacción a errores, sonidos o bloqueos]].
            \item \textbf{Valores:} [[Sentimiento de logro, orgullo por autonomía]].
        },
        learningStyle = {
            \item \textbf{Perfil de aprendizaje:} [[\textbf{Visual} / \textbf{Auditivo} / \textbf{Kinestésico}]].
            \item \textbf{Canales óptimos:} [[Fotos reales / Pictogramas ARASAAC / Vídeos]].
            \item \textbf{Metodología:} [[Modelado guiado, repetición, ensayo y éxito]].
        },
        criticality = {
            \item \textbf{Consecuencias de un error:} \textit{Emocional:} [[¿Bloqueo, crisis de frustración, rechazo?]] \textit{Operativo:} [[¿Desorientación, interrupción de la tarea?]]
            \item \textbf{Tolerancia del software:} [[Acciones no destructivas, confirmación siempre asistida, botón SOS visible]].
        },
        escenario = {[[Narrar en 3-4 líneas una situación real: a qué hora usa la agenda, cómo anticipa su rutina con apoyos visuales y cómo supera el momento crítico logrando autonomía sin estrés.]]},
        badgesAjustes = {
            \badge{primaryColor}{\faKey\ Login: [[Método]]} \hspace{2mm}
            \badge{secondaryColor}{\faIcons\ [[ARASAAC / Fotos]]} \hspace{2mm}
            \badge{accentColor}{\faClock\ [[Temporizador]]} \hspace{2mm}
            \badge{purpleColor}{\faVolumeUp\ [[Voz TTS]]}%
        }
    },
    bannerTitulo/.store in = \fAlBannerTitulo,
    subtitulo/.store in = \fAlSubtitulo,
    equipo/.store in = \fAlEquipo,
    nombre/.store in = \fAlNombre,
    nombreCorto/.store in = \fAlNombreCorto,
    edad/.store in = \fAlEdad,
    icono/.store in = \fAlIcono,
    foto/.store in = \fAlFoto,
    badgeRol/.store in = \fAlBadgeRol,
    badgeTarea/.store in = \fAlBadgeTarea,
    centro/.store in = \fAlCentro,
    cita/.store in = \fAlCita,
    demographic/.store in = \fAlDemographic,
    computerExp/.store in = \fAlComputerExp,
    productExp/.store in = \fAlProductExp,
    tasks/.store in = \fAlTasks,
    domainKnowledge/.store in = \fAlDomainKnowledge,
    technology/.store in = \fAlTechnology,
    attitudes/.store in = \fAlAttitudes,
    learningStyle/.store in = \fAlLearningStyle,
    criticality/.store in = \fAlCriticality,
    escenario/.store in = \fAlEscenario,
    badgesAjustes/.store in = \fAlBadgesAjustes,
}

% COMANDO PRINCIPAL \fichaAlumno
\newcommand{\fichaAlumno}[1][]{%
    \begingroup
    \pgfkeys{/fichaAlumno, default, #1}%
    % 1. BANNER SUPERIOR
    \begin{tcolorbox}[
        enhanced,
        colback=primaryColor,
        colframe=primaryColor,
        arc=1.6mm,
        left=7pt, right=7pt, top=2.5pt, bottom=2.5pt,
        boxsep=0pt, coltext=white
    ]
        \begin{minipage}[c]{0.70\textwidth}
            {\normalsize\textbf{\fAlBannerTitulo}}\\[0.2mm]
            {\tiny \fAlSubtitulo}
        \end{minipage}\hfill
        \begin{minipage}[c]{0.29\textwidth}
            \raggedleft\tiny
            \fAlEquipo
        \end{minipage}
    \end{tcolorbox}
    \vspace{0.8mm}
    % 2. TARJETA DE IDENTIDAD DEL ALUMNO
    \begin{tcolorbox}[
        enhanced,
        colback=white,
        colframe=primaryColor!65,
        arc=2.0mm,
        boxrule=0.8pt,
        left=6pt, right=6pt, top=3.0pt, bottom=3.0pt,
        boxsep=0pt,
        drop shadow=black!5
    ]
        \begin{minipage}[c]{0.10\textwidth}
            \centering
            \begin{tikzpicture}
                \ifx\fAlFoto\empty
                    \draw[fill=primaryColor!8, draw=primaryColor!30, line width=0.5pt, rounded corners=2mm] (0,0) rectangle (1.4, 1.55);
                    \node at (0.7, 0.95) {\color{primaryColor}\scalebox{1.5}{\fAlIcono}};
                    \node[align=center] at (0.7, 0.28) {\tiny\textbf{\color{primaryColor}\fAlNombreCorto}\\\fontsize{4pt}{5pt}\selectfont\color{textMuted}\fAlEdad};
                \else
                    \begin{scope}
                        \clip[rounded corners=2mm] (0,0) rectangle (1.45, 1.55);
                        \node[inner sep=0pt] at (0.725, 0.775) {\includegraphics[width=1.45cm,height=1.55cm,keepaspectratio=false]{\fAlFoto}};
                    \end{scope}
                    \draw[draw=primaryColor!50, line width=0.8pt, rounded corners=2mm] (0,0) rectangle (1.45, 1.55);
                \fi
            \end{tikzpicture}
        \end{minipage}\hfill
        \begin{minipage}[c]{0.88\textwidth}
            \begin{minipage}[b]{0.68\textwidth}\raggedright
                {\large\textbf{\color{primaryColor}\fAlNombre}} \hspace{2mm}
                \badge{secondaryColor}{\fAlBadgeRol} \badge{accentColor}{\fAlBadgeTarea}
            \end{minipage}\hfill
            \begin{minipage}[b]{0.31\textwidth}
                \raggedleft\small\textbf{\color{primaryColor}\faSchool\ \fAlCentro}
            \end{minipage}\\[0.6mm]
            \tikz[baseline] \node[fill=primaryColor!5, draw=primaryColor!25, rounded corners=1.2mm, inner sep=2.5pt, text width=0.98\linewidth] {
                \small\itshape\color{primaryColor}
                \faQuoteLeft\ ``\fAlCita''\ \faQuoteRight
            };
        \end{minipage}
    \end{tcolorbox}
    \vspace{0.8mm}
    % 3. LAS 9 CAJAS DE ALUMNOS (4 FILAS X 2 COLS + 1 FILA COMPLETA)
    \noindent\begin{minipage}[t]{0.485\textwidth}
        \begin{tcolorbox}[cardAlumnoBox, colframe=primaryColor!50, coltitle=primaryColor, title={\faIdCard\ 1. Demographic characteristics}]
            \fontsize{8.2pt}{9.8pt}\selectfont
            \begin{itemize}
                \fAlDemographic
            \end{itemize}
        \end{tcolorbox}
    \end{minipage}\hfill
    \begin{minipage}[t]{0.485\textwidth}
        \begin{tcolorbox}[cardAlumnoBox, colframe=secondaryColor!60, coltitle=secondaryColor!90!black, title={\faLaptop\ 2. Computer experience}]
            \fontsize{8.2pt}{9.8pt}\selectfont
            \begin{itemize}
                \fAlComputerExp
            \end{itemize}
        \end{tcolorbox}
    \end{minipage}
    \vspace{0.8mm}
    \noindent\begin{minipage}[t]{0.485\textwidth}
        \begin{tcolorbox}[cardAlumnoBox, colframe=secondaryColor!60, coltitle=secondaryColor!90!black, title={\faSearchPlus\ 3. Specific product experience}]
            \fontsize{8.2pt}{9.8pt}\selectfont
            \begin{itemize}
                \fAlProductExp
            \end{itemize}
        \end{tcolorbox}
    \end{minipage}\hfill
    \begin{minipage}[t]{0.485\textwidth}
        \begin{tcolorbox}[cardAlumnoBox, colframe=greenColor!65, coltitle=greenColor!90!black, title={\faTasks\ 4. Tasks (Primary and secondary)}]
            \fontsize{8.2pt}{9.8pt}\selectfont
            \begin{itemize}
                \fAlTasks
            \end{itemize}
        \end{tcolorbox}
    \end{minipage}
    \vspace{0.8mm}
    \noindent\begin{minipage}[t]{0.485\textwidth}
        \begin{tcolorbox}[cardAlumnoBox, colframe=accentColor!65, coltitle=accentColor!90!black, title={\faBrain\ 5. Domain knowledge}]
            \fontsize{8.2pt}{9.8pt}\selectfont
            \begin{itemize}
                \fAlDomainKnowledge
            \end{itemize}
        \end{tcolorbox}
    \end{minipage}\hfill
    \begin{minipage}[t]{0.485\textwidth}
        \begin{tcolorbox}[cardAlumnoBox, colframe=tealColor!60, coltitle=tealColor!90!black, title={\faTools\ 6. Technology available}]
            \fontsize{8.2pt}{9.8pt}\selectfont
            \begin{itemize}
                \fAlTechnology
            \end{itemize}
        \end{tcolorbox}
    \end{minipage}
    \vspace{0.8mm}
    \noindent\begin{minipage}[t]{0.485\textwidth}
        \begin{tcolorbox}[cardAlumnoBox, colframe=purpleColor!60, coltitle=purpleColor!90!black, title={\faHeart\ 7. Attitudes and values}]
            \fontsize{8.2pt}{9.8pt}\selectfont
            \begin{itemize}
                \fAlAttitudes
            \end{itemize}
        \end{tcolorbox}
    \end{minipage}\hfill
    \begin{minipage}[t]{0.485\textwidth}
        \begin{tcolorbox}[cardAlumnoBox, colframe=purpleColor!60, coltitle=purpleColor!90!black, title={\faEye\ 8. Learning style}]
            \fontsize{8.2pt}{9.8pt}\selectfont
            \begin{itemize}
                \fAlLearningStyle
            \end{itemize}
        \end{tcolorbox}
    \end{minipage}
    \vspace{0.8mm}
    % FILA 5: Punto 9 (ANCHO COMPLETO — CRITICALITY OF ERRORS)
    \noindent\begin{minipage}[t]{\textwidth}
        \begin{tcolorbox}[cardAlumnoWideBox, colframe=dangerColor!65, coltitle=dangerColor!90!black, title={\faExclamationTriangle\ 9. Criticality of errors (Consecuencias de un error y tolerancia del software)}]
            \fontsize{8.2pt}{9.8pt}\selectfont
            \begin{itemize}
                \fAlCriticality
            \end{itemize}
        \end{tcolorbox}
    \end{minipage}
    \vspace{0.8mm}
    % 4. BARRA INFERIOR DE INTEGRACIÓN CeRA
    \begin{tcolorbox}[
        enhanced,
        colback=primaryColor!4,
        colframe=primaryColor!45,
        arc=1.6mm,
        left=6pt, right=6pt, top=2.5pt, bottom=2.5pt,
        boxsep=0pt
    ]
        \scriptsize
        \textbf{\color{primaryColor}\faMapSigns\ Escenario de Uso Contextual (Día típico):} \fAlEscenario\\
        \vspace{0.3mm}
        \hrule height 0.35pt \relax
        \vspace{0.5mm}
        \centerline{%
            \textbf{\color{primaryColor}\faSlidersH\ Ajustes UI Accesible (CeRA):} \hspace{2mm}
            \fAlBadgesAjustes%
        }
    \end{tcolorbox}
    \newpage
    \endgroup
}

% COMANDO PARA PLANTILLA EN BLANCO DE ALUMNADO
\newcommand{\fichaAlumnoBlanco}{%
    \fichaAlumno[%
        bannerTitulo={\faUserEdit\ FICHA DE PERSONA: ALUMNADO \textbar\ PLANTILLA DE TRABAJO},
        cita={[Cita textual que exprese cómo vive su rutina, qué le ayuda a anticipar y qué le genera ansiedad o rechazo].}%
    ]%
}

\ifdefined\isSubFile
    \expandafter\endinput
\fi

\begin{document}
\fichaAlumnoBlanco
\end{document}
 y luego invocar \fichaAlumno[...] o \fichaAlumnoBlanco
% ==============================================================================

\ifdefined\isSubFile
    % Si se incluye desde otro documento, los paquetes ya están cargados
\else
    \documentclass[10pt]{article}
    \usepackage[utf8]{inputenc}
    \usepackage[spanish,es-tabla]{babel}
    \usepackage[margin=0.8cm,top=0.6cm,bottom=0.6cm]{geometry}
    \usepackage{helvet}
    \renewcommand{\familydefault}{\sfdefault}
    \usepackage{microtype}
    \usepackage{graphicx}
    \usepackage{xcolor}
    \usepackage{tikz}
    \usepackage{enumitem}
    \usepackage{fontawesome5}
    \usepackage{tcolorbox}
    \tcbuselibrary{skins}
    \pagestyle{empty}
\fi

% Colores del sistema
\providecolor{primaryColor}{RGB}{20, 65, 105}     % Azul corporativo
\providecolor{secondaryColor}{RGB}{35, 115, 170}  % Azul medio accesible
\providecolor{tealColor}{RGB}{20, 130, 140}       % Verde azulado
\providecolor{greenColor}{RGB}{30, 135, 75}       % Verde bosque
\providecolor{accentColor}{RGB}{215, 105, 25}     % Naranja cálido
\providecolor{purpleColor}{RGB}{120, 60, 150}     % Púrpura suave
\providecolor{dangerColor}{RGB}{185, 45, 35}      % Rojo criticidad
\providecolor{bgCard}{RGB}{252, 253, 255}         % Fondo limpio
\providecolor{borderCard}{RGB}{215, 225, 235}     % Borde sutil
\providecolor{textDark}{RGB}{30, 40, 50}          % Texto principal oscuro
\providecolor{textMuted}{RGB}{95, 110, 125}       % Texto secundario

% Configuración de listas
\setlist[itemize]{leftmargin=3.0mm, itemsep=0.4pt, parsep=0pt, topsep=0.4pt, partopsep=0pt}

% Estilos de caja fija para Alumno
\tcbset{
    cardAlumnoBox/.style={
        enhanced,
        boxrule=0.65pt,
        arc=1.8mm,
        height=32.0mm,
        left=5.0pt, right=5.0pt, top=1.8pt, bottom=1.8pt,
        boxsep=0pt,
        colback=bgCard,
        colframe=borderCard,
        fonttitle=\bfseries\small,
        fontupper=\fontsize{8.5pt}{10.0pt}\selectfont,
    },
    cardAlumnoWideBox/.style={
        enhanced,
        boxrule=0.65pt,
        arc=1.8mm,
        height=22.5mm,
        left=6.0pt, right=6.0pt, top=2.0pt, bottom=2.0pt,
        boxsep=0pt,
        colback=bgCard,
        colframe=borderCard,
        fonttitle=\bfseries\small,
    }
}

% Macro para badges visuales (si no existe)
\providecommand{\badge}[2]{%
    \tikz[baseline=(b.base)] \node[fill=#1!12, text=#1!90!black, draw=#1!35, line width=0.4pt, rounded corners=1.2mm, inner sep=2.0pt, font=\bfseries\scriptsize] (b) {#2};%
}

% Definición de claves con pgfkeys para \fichaAlumno
\pgfkeys{
    /fichaAlumno/.is family, /fichaAlumno,
    default/.style = {
        bannerTitulo = {\faUserGraduate\ FICHA DE PERSONA: ALUMNADO (METODOLOGÍA DCU)},
        subtitulo = {Proyecto AUTOVIDA (26-175) \textbar\ Reto CeRA — Agenda Accesible PTVAL \textbar\ Asignatura MDA (UGR)},
        equipo = {\textbf{Equipo CeRA-MDA} (5 integrantes)\\\ Course 2026/2027 \textbar\ \textbf{Práctica 1}},
        nombre = {[Nombre y Apellidos del Alumno]},
        nombreCorto = {[ALUMNO]},
        edad = {[Edad]},
        icono = {\faUserGraduate},
        foto = {},
        badgeRol = {[Perfil NEAE / Aula]},
        badgeTarea = {Tareas Rotativas},
        centro = {CEE Purísima Concepción},
        cita = {[Cita textual representativa que exprese cómo vive su rutina, qué le ayuda a anticipar y qué le genera ansiedad o rechazo].},
        demographic = {
            \item \textbf{Edad y Género (Age / Gender):} [[Edad] años, [Género]].
            \item \textbf{Ubicación (Location):} [[Granada capital / Municipio]].
            \item \textbf{Perfil diagnóstico:} [[Diagnóstico cognitivo, motor o sensorial]].
        },
        computerExp = {
            \item \textbf{Habilidades informáticas:} [[Nivel táctil: autónomo, asistido]].
            \item \textbf{Años de experiencia:} [[Años de contacto con tablets]].
            \item \textbf{Hábitos de uso:} [[Apps frecuentes: vídeos, pictogramas, juegos]].
            \item \textbf{Autonomía digital:} [[Baja / Media / Asistida]].
        },
        productExp = {
            \item \textbf{Productos similares:} [[PECS en papel, velcro, apps previas]].
            \item \textbf{Limitaciones actuales:} [[Pérdida de material, rigidez, frustración]].
            \item \textbf{Tendencias de uso:} [[Receptividad a imágenes reales vs pictogramas]].
        },
        tasks = {
            \item \textbf{Tareas Primarias:} 1) [[Login accesible]] $\rightarrow$ 2) [[Ver tarea actual]] $\rightarrow$ 3) [[Ejecutar pasos guiados]] $\rightarrow$ 4) [[Confirmar finalización]].
            \item \textbf{Tareas Secundarias:} [[Ver siguiente tarea / pedir ayuda]].
        },
        domainKnowledge = {
            \item \textbf{Comprensión del área:} [[Orientación en el colegio, aulas y profesores]].
            \item \textbf{Comprensión temporal:} [[Visual (disco/reloj de arena), no numérica]].
            \item \textbf{Destreza procedimental:} [[Grado de autonomía en la tarea asignada]].
        },
        technology = {
            \item \textbf{Hardware:} [[Tablet Android / pulsador externo]].
            \item \textbf{Software:} [[Android / síntesis de voz]].
            \item \textbf{Herramientas habituales:} [[Conexión Wi-Fi / atril]].
        },
        attitudes = {
            \item \textbf{Preferencias:} [[Refuerzos sensoriales favoritos, limpieza visual]].
            \item \textbf{Miedo a la tecnología:} [[Reacción a errores, sonidos o bloqueos]].
            \item \textbf{Valores:} [[Sentimiento de logro, orgullo por autonomía]].
        },
        learningStyle = {
            \item \textbf{Perfil de aprendizaje:} [[\textbf{Visual} / \textbf{Auditivo} / \textbf{Kinestésico}]].
            \item \textbf{Canales óptimos:} [[Fotos reales / Pictogramas ARASAAC / Vídeos]].
            \item \textbf{Metodología:} [[Modelado guiado, repetición, ensayo y éxito]].
        },
        criticality = {
            \item \textbf{Consecuencias de un error:} \textit{Emocional:} [[¿Bloqueo, crisis de frustración, rechazo?]] \textit{Operativo:} [[¿Desorientación, interrupción de la tarea?]]
            \item \textbf{Tolerancia del software:} [[Acciones no destructivas, confirmación siempre asistida, botón SOS visible]].
        },
        escenario = {[[Narrar en 3-4 líneas una situación real: a qué hora usa la agenda, cómo anticipa su rutina con apoyos visuales y cómo supera el momento crítico logrando autonomía sin estrés.]]},
        badgesAjustes = {
            \badge{primaryColor}{\faKey\ Login: [[Método]]} \hspace{2mm}
            \badge{secondaryColor}{\faIcons\ [[ARASAAC / Fotos]]} \hspace{2mm}
            \badge{accentColor}{\faClock\ [[Temporizador]]} \hspace{2mm}
            \badge{purpleColor}{\faVolumeUp\ [[Voz TTS]]}%
        }
    },
    bannerTitulo/.store in = \fAlBannerTitulo,
    subtitulo/.store in = \fAlSubtitulo,
    equipo/.store in = \fAlEquipo,
    nombre/.store in = \fAlNombre,
    nombreCorto/.store in = \fAlNombreCorto,
    edad/.store in = \fAlEdad,
    icono/.store in = \fAlIcono,
    foto/.store in = \fAlFoto,
    badgeRol/.store in = \fAlBadgeRol,
    badgeTarea/.store in = \fAlBadgeTarea,
    centro/.store in = \fAlCentro,
    cita/.store in = \fAlCita,
    demographic/.store in = \fAlDemographic,
    computerExp/.store in = \fAlComputerExp,
    productExp/.store in = \fAlProductExp,
    tasks/.store in = \fAlTasks,
    domainKnowledge/.store in = \fAlDomainKnowledge,
    technology/.store in = \fAlTechnology,
    attitudes/.store in = \fAlAttitudes,
    learningStyle/.store in = \fAlLearningStyle,
    criticality/.store in = \fAlCriticality,
    escenario/.store in = \fAlEscenario,
    badgesAjustes/.store in = \fAlBadgesAjustes,
}

% COMANDO PRINCIPAL \fichaAlumno
\newcommand{\fichaAlumno}[1][]{%
    \begingroup
    \pgfkeys{/fichaAlumno, default, #1}%
    % 1. BANNER SUPERIOR
    \begin{tcolorbox}[
        enhanced,
        colback=primaryColor,
        colframe=primaryColor,
        arc=1.6mm,
        left=7pt, right=7pt, top=2.5pt, bottom=2.5pt,
        boxsep=0pt, coltext=white
    ]
        \begin{minipage}[c]{0.70\textwidth}
            {\normalsize\textbf{\fAlBannerTitulo}}\\[0.2mm]
            {\tiny \fAlSubtitulo}
        \end{minipage}\hfill
        \begin{minipage}[c]{0.29\textwidth}
            \raggedleft\tiny
            \fAlEquipo
        \end{minipage}
    \end{tcolorbox}
    \vspace{0.8mm}
    % 2. TARJETA DE IDENTIDAD DEL ALUMNO
    \begin{tcolorbox}[
        enhanced,
        colback=white,
        colframe=primaryColor!65,
        arc=2.0mm,
        boxrule=0.8pt,
        left=6pt, right=6pt, top=3.0pt, bottom=3.0pt,
        boxsep=0pt,
        drop shadow=black!5
    ]
        \begin{minipage}[c]{0.10\textwidth}
            \centering
            \begin{tikzpicture}
                \ifx\fAlFoto\empty
                    \draw[fill=primaryColor!8, draw=primaryColor!30, line width=0.5pt, rounded corners=2mm] (0,0) rectangle (1.4, 1.55);
                    \node at (0.7, 0.95) {\color{primaryColor}\scalebox{1.5}{\fAlIcono}};
                    \node[align=center] at (0.7, 0.28) {\tiny\textbf{\color{primaryColor}\fAlNombreCorto}\\\fontsize{4pt}{5pt}\selectfont\color{textMuted}\fAlEdad};
                \else
                    \begin{scope}
                        \clip[rounded corners=2mm] (0,0) rectangle (1.45, 1.55);
                        \node[inner sep=0pt] at (0.725, 0.775) {\includegraphics[width=1.45cm,height=1.55cm,keepaspectratio=false]{\fAlFoto}};
                    \end{scope}
                    \draw[draw=primaryColor!50, line width=0.8pt, rounded corners=2mm] (0,0) rectangle (1.45, 1.55);
                \fi
            \end{tikzpicture}
        \end{minipage}\hfill
        \begin{minipage}[c]{0.88\textwidth}
            \begin{minipage}[b]{0.68\textwidth}\raggedright
                {\large\textbf{\color{primaryColor}\fAlNombre}} \hspace{2mm}
                \badge{secondaryColor}{\fAlBadgeRol} \badge{accentColor}{\fAlBadgeTarea}
            \end{minipage}\hfill
            \begin{minipage}[b]{0.31\textwidth}
                \raggedleft\small\textbf{\color{primaryColor}\faSchool\ \fAlCentro}
            \end{minipage}\\[0.6mm]
            \tikz[baseline] \node[fill=primaryColor!5, draw=primaryColor!25, rounded corners=1.2mm, inner sep=2.5pt, text width=0.98\linewidth] {
                \small\itshape\color{primaryColor}
                \faQuoteLeft\ ``\fAlCita''\ \faQuoteRight
            };
        \end{minipage}
    \end{tcolorbox}
    \vspace{0.8mm}
    % 3. LAS 9 CAJAS DE ALUMNOS (4 FILAS X 2 COLS + 1 FILA COMPLETA)
    \noindent\begin{minipage}[t]{0.485\textwidth}
        \begin{tcolorbox}[cardAlumnoBox, colframe=primaryColor!50, coltitle=primaryColor, title={\faIdCard\ 1. Demographic characteristics}]
            \fontsize{8.2pt}{9.8pt}\selectfont
            \begin{itemize}
                \fAlDemographic
            \end{itemize}
        \end{tcolorbox}
    \end{minipage}\hfill
    \begin{minipage}[t]{0.485\textwidth}
        \begin{tcolorbox}[cardAlumnoBox, colframe=secondaryColor!60, coltitle=secondaryColor!90!black, title={\faLaptop\ 2. Computer experience}]
            \fontsize{8.2pt}{9.8pt}\selectfont
            \begin{itemize}
                \fAlComputerExp
            \end{itemize}
        \end{tcolorbox}
    \end{minipage}
    \vspace{0.8mm}
    \noindent\begin{minipage}[t]{0.485\textwidth}
        \begin{tcolorbox}[cardAlumnoBox, colframe=secondaryColor!60, coltitle=secondaryColor!90!black, title={\faSearchPlus\ 3. Specific product experience}]
            \fontsize{8.2pt}{9.8pt}\selectfont
            \begin{itemize}
                \fAlProductExp
            \end{itemize}
        \end{tcolorbox}
    \end{minipage}\hfill
    \begin{minipage}[t]{0.485\textwidth}
        \begin{tcolorbox}[cardAlumnoBox, colframe=greenColor!65, coltitle=greenColor!90!black, title={\faTasks\ 4. Tasks (Primary and secondary)}]
            \fontsize{8.2pt}{9.8pt}\selectfont
            \begin{itemize}
                \fAlTasks
            \end{itemize}
        \end{tcolorbox}
    \end{minipage}
    \vspace{0.8mm}
    \noindent\begin{minipage}[t]{0.485\textwidth}
        \begin{tcolorbox}[cardAlumnoBox, colframe=accentColor!65, coltitle=accentColor!90!black, title={\faBrain\ 5. Domain knowledge}]
            \fontsize{8.2pt}{9.8pt}\selectfont
            \begin{itemize}
                \fAlDomainKnowledge
            \end{itemize}
        \end{tcolorbox}
    \end{minipage}\hfill
    \begin{minipage}[t]{0.485\textwidth}
        \begin{tcolorbox}[cardAlumnoBox, colframe=tealColor!60, coltitle=tealColor!90!black, title={\faTools\ 6. Technology available}]
            \fontsize{8.2pt}{9.8pt}\selectfont
            \begin{itemize}
                \fAlTechnology
            \end{itemize}
        \end{tcolorbox}
    \end{minipage}
    \vspace{0.8mm}
    \noindent\begin{minipage}[t]{0.485\textwidth}
        \begin{tcolorbox}[cardAlumnoBox, colframe=purpleColor!60, coltitle=purpleColor!90!black, title={\faHeart\ 7. Attitudes and values}]
            \fontsize{8.2pt}{9.8pt}\selectfont
            \begin{itemize}
                \fAlAttitudes
            \end{itemize}
        \end{tcolorbox}
    \end{minipage}\hfill
    \begin{minipage}[t]{0.485\textwidth}
        \begin{tcolorbox}[cardAlumnoBox, colframe=purpleColor!60, coltitle=purpleColor!90!black, title={\faEye\ 8. Learning style}]
            \fontsize{8.2pt}{9.8pt}\selectfont
            \begin{itemize}
                \fAlLearningStyle
            \end{itemize}
        \end{tcolorbox}
    \end{minipage}
    \vspace{0.8mm}
    % FILA 5: Punto 9 (ANCHO COMPLETO — CRITICALITY OF ERRORS)
    \noindent\begin{minipage}[t]{\textwidth}
        \begin{tcolorbox}[cardAlumnoWideBox, colframe=dangerColor!65, coltitle=dangerColor!90!black, title={\faExclamationTriangle\ 9. Criticality of errors (Consecuencias de un error y tolerancia del software)}]
            \fontsize{8.2pt}{9.8pt}\selectfont
            \begin{itemize}
                \fAlCriticality
            \end{itemize}
        \end{tcolorbox}
    \end{minipage}
    \vspace{0.8mm}
    % 4. BARRA INFERIOR DE INTEGRACIÓN CeRA
    \begin{tcolorbox}[
        enhanced,
        colback=primaryColor!4,
        colframe=primaryColor!45,
        arc=1.6mm,
        left=6pt, right=6pt, top=2.5pt, bottom=2.5pt,
        boxsep=0pt
    ]
        \scriptsize
        \textbf{\color{primaryColor}\faMapSigns\ Escenario de Uso Contextual (Día típico):} \fAlEscenario\\
        \vspace{0.3mm}
        \hrule height 0.35pt \relax
        \vspace{0.5mm}
        \centerline{%
            \textbf{\color{primaryColor}\faSlidersH\ Ajustes UI Accesible (CeRA):} \hspace{2mm}
            \fAlBadgesAjustes%
        }
    \end{tcolorbox}
    \newpage
    \endgroup
}

% COMANDO PARA PLANTILLA EN BLANCO DE ALUMNADO
\newcommand{\fichaAlumnoBlanco}{%
    \fichaAlumno[%
        bannerTitulo={\faUserEdit\ FICHA DE PERSONA: ALUMNADO \textbar\ PLANTILLA DE TRABAJO},
        cita={[Cita textual que exprese cómo vive su rutina, qué le ayuda a anticipar y qué le genera ansiedad o rechazo].}%
    ]%
}

\ifdefined\isSubFile
    \expandafter\endinput
\fi

\begin{document}
\fichaAlumnoBlanco
\end{document}

% ==============================================================================
% PLANTILLA_ADULTO.TEX — MACRO Y DEFINICIÓN DE FICHA PARA TUTORES Y FAMILIAS
% Asignatura: Metodologías de Desarrollo Ágil (MDA) — Universidad de Granada
% Proyecto AUTOVIDA (26-175) | Reto CeRA — Agenda Accesible PTVAL
% ==============================================================================
% Campos canónicos: 11 (Con Occupation y Education, sin Company Information)
% Uso: % ==============================================================================
% PLANTILLA_ADULTO.TEX — MACRO Y DEFINICIÓN DE FICHA PARA TUTORES Y FAMILIAS
% Asignatura: Metodologías de Desarrollo Ágil (MDA) — Universidad de Granada
% Proyecto AUTOVIDA (26-175) | Reto CeRA — Agenda Accesible PTVAL
% ==============================================================================
% Campos canónicos: 11 (Con Occupation y Education, sin Company Information)
% Uso: % ==============================================================================
% PLANTILLA_ADULTO.TEX — MACRO Y DEFINICIÓN DE FICHA PARA TUTORES Y FAMILIAS
% Asignatura: Metodologías de Desarrollo Ágil (MDA) — Universidad de Granada
% Proyecto AUTOVIDA (26-175) | Reto CeRA — Agenda Accesible PTVAL
% ==============================================================================
% Campos canónicos: 11 (Con Occupation y Education, sin Company Information)
% Uso: \input{plantilla_adulto.tex} y luego invocar \fichaAdulto[...] o \fichaAdultoBlanco
% ==============================================================================

\ifdefined\isSubFile
    % Si se incluye desde otro documento, los paquetes ya están cargados
\else
    \documentclass[10pt]{article}
    \usepackage[utf8]{inputenc}
    \usepackage[spanish,es-tabla]{babel}
    \usepackage[margin=0.8cm,top=0.6cm,bottom=0.6cm]{geometry}
    \usepackage{helvet}
    \renewcommand{\familydefault}{\sfdefault}
    \usepackage{microtype}
    \usepackage{graphicx}
    \usepackage{xcolor}
    \usepackage{tikz}
    \usepackage{enumitem}
    \usepackage{fontawesome5}
    \usepackage{tcolorbox}
    \tcbuselibrary{skins}
    \pagestyle{empty}
\fi

% Colores del sistema
\providecolor{primaryColor}{RGB}{20, 65, 105}     % Azul corporativo
\providecolor{secondaryColor}{RGB}{35, 115, 170}  % Azul medio accesible
\providecolor{tealColor}{RGB}{20, 130, 140}       % Verde azulado
\providecolor{greenColor}{RGB}{30, 135, 75}       % Verde bosque
\providecolor{accentColor}{RGB}{215, 105, 25}     % Naranja cálido
\providecolor{purpleColor}{RGB}{120, 60, 150}     % Púrpura suave
\providecolor{dangerColor}{RGB}{185, 45, 35}      % Rojo criticidad
\providecolor{bgCard}{RGB}{252, 253, 255}         % Fondo limpio
\providecolor{borderCard}{RGB}{215, 225, 235}     % Borde sutil
\providecolor{textDark}{RGB}{30, 40, 50}          % Texto principal oscuro
\providecolor{textMuted}{RGB}{95, 110, 125}       % Texto secundario

% Configuración de listas
\setlist[itemize]{leftmargin=3.0mm, itemsep=0.4pt, parsep=0pt, topsep=0.4pt, partopsep=0pt}

% Estilos de caja fija para Adulto
\tcbset{
    cardAdultoBox/.style={
        enhanced,
        boxrule=0.65pt,
        arc=1.8mm,
        height=27.2mm,
        left=5.0pt, right=5.0pt, top=1.5pt, bottom=1.5pt,
        boxsep=0pt,
        colback=bgCard,
        colframe=borderCard,
        fonttitle=\bfseries\small,
        fontupper=\fontsize{8.2pt}{9.6pt}\selectfont,
    },
    cardAdultoWideBox/.style={
        enhanced,
        boxrule=0.65pt,
        arc=1.8mm,
        height=20.0mm,
        left=6.0pt, right=6.0pt, top=2.0pt, bottom=2.0pt,
        boxsep=0pt,
        colback=bgCard,
        colframe=borderCard,
        fonttitle=\bfseries\small,
    }
}

% Macro para badges visuales (si no existe)
\providecommand{\badge}[2]{%
    \tikz[baseline=(b.base)] \node[fill=#1!12, text=#1!90!black, draw=#1!35, line width=0.4pt, rounded corners=1.2mm, inner sep=2.0pt, font=\bfseries\scriptsize] (b) {#2};%
}

% Definición de claves con pgfkeys para \fichaAdulto
\pgfkeys{
    /fichaAdulto/.is family, /fichaAdulto,
    default/.style = {
        bannerTitulo = {\faChalkboardTeacher\ FICHA DE PERSONA: TUTORES Y FAMILIAS (METODOLOGÍA DCU)},
        subtitulo = {Proyecto AUTOVIDA (26-175) \textbar\ Reto CeRA — Agenda Accesible PTVAL \textbar\ Asignatura MDA (UGR)},
        equipo = {\textbf{Equipo CeRA-MDA} (5 integrantes)\\\ Course 2026/2027 \textbar\ \textbf{Práctica 1}},
        nombre = {[Nombre y Apellidos del Usuario]},
        nombreCorto = {[ADULTO]},
        edad = {[Edad]},
        icono = {\faChalkboardTeacher},
        foto = {},
        badgeRol = {[Rol: Tutor / Familiar / Admin]},
        badgeTarea = {[Ámbito de Actuación]},
        centro = {CEE Purísima Concepción},
        cita = {[Cita textual representativa que exprese sus objetivos, inquietudes de gestión o necesidades respecto a la app CeRA].},
        demographic = {
            \item \textbf{Edad y Género:} [[Edad] años, [Género]].
            \item \textbf{Ubicación:} [[Ciudad / Zona residencial]].
        },
        occupation = {
            \item \textbf{Puesto actual:} [[Puesto docente, profesión o rol familiar]].
            \item \textbf{Antigüedad:} [[Años en el centro/entidad] \textbar\ [Años en puesto]].
            \item \textbf{Responsabilidades:} [[Supervisión, docencia, apoyo en el hogar]].
        },
        education = {
            \item \textbf{Titulación (Degree):} [[Nivel formativo o titulación alcanzada]].
            \item \textbf{Especialidad (Major):} [[Educación Especial, Pedagogía, etc.]].
            \item \textbf{Cursos realizados:} [[Formación en NEAE, SAAC, TIC o gestión]].
        },
        computerExp = {
            \item \textbf{Habilidades informáticas:} [[Nivel de competencia digital: medio/avanzado]].
            \item \textbf{Años de experiencia:} [[Años de uso regular de dispositivos]].
            \item \textbf{Hábitos de uso:} [[Navegación web, apps de gestión, smartphones]].
        },
        productExp = {
            \item \textbf{Productos similares:} [[Otras herramientas usadas, papel, agendas]].
            \item \textbf{Limitaciones actuales:} [[Sobrecarga de trabajo, falta de personalización]].
            \item \textbf{Expectativas:} [[Ahorro de tiempo y canal de dudas]].
        },
        tasks = {
            \item \textbf{Tareas Primarias:} 1) [[Configurar tareas]] $\rightarrow$ 2) [[Asignar a alumnos]] $\rightarrow$ 3) [[Supervisar ejecución]] $\rightarrow$ 4) [[Gestionar comedor/avisos]].
            \item \textbf{Tareas Secundarias:} [[Revisar estadísticas / chat con familias]].
        },
        domainKnowledge = {
            \item \textbf{Comprensión del área:} [[Conocimiento sobre necesidades NEAE y rutinas]].
            \item \textbf{Conocimiento del usuario:} [[Grado de personalización que requiere cada estudiante]].
        },
        technology = {
            \item \textbf{Hardware:} [[Ordenador portátil / tablet / móvil Android]].
            \item \textbf{Software:} [[Navegador web / app móvil]].
            \item \textbf{Herramientas:} [[Conexión a internet]].
        },
        attitudes = {
            \item \textbf{Preferencias:} [[Sencillez, rapidez operativa y fiabilidad en tiempo real]].
            \item \textbf{Miedo a la tecnología:} [[Preocupación por pérdida de datos o fallos de red]].
            \item \textbf{Valores:} [[Foco en el progreso y bienestar del alumnado]].
        },
        learningStyle = {
            \item \textbf{Perfil de aprendizaje:} [[\textbf{Práctico} / \textbf{Estructurado} / \textbf{Visual}]].
            \item \textbf{Canales óptimos:} [[Tutoriales cortos, interfaz intuitiva sin manual]].
            \item \textbf{Curva de adopción:} [[Inmediata si soluciona problemas reales]].
        },
        criticality = {
            \item \textbf{Consecuencias de un error:} \textit{Logístico:} [[¿Dietas erróneas, tareas desincronizadas?]] \textit{Pedagógico:} [[¿Confusión en el horario del alumno, estrés?]]
            \item \textbf{Tolerancia del software:} [[Alertas claras ante cambios sensibles, confirmación de guardado]].
        },
        escenario = {[[Describir brevemente en 3-4 líneas cómo accede al panel de la app, planifica o revisa las tareas de sus estudiantes y cómo la herramienta le facilita la jornada eliminando fricciones organizativas.]]},
        badgesAjustes = {
            \badge{tealColor}{\faDesktop\ [[Panel Web]]} \hspace{2mm}
            \badge{primaryColor}{\faCalendarCheck\ [[Planificación]]} \hspace{2mm}
            \badge{accentColor}{\faUtensils\ [[Comedor/Dietas]]} \hspace{2mm}
            \badge{greenColor}{\faCheckCircle\ [[Supervisión]]}%
        }
    },
    bannerTitulo/.store in = \fAdBannerTitulo,
    subtitulo/.store in = \fAdSubtitulo,
    equipo/.store in = \fAdEquipo,
    nombre/.store in = \fAdNombre,
    nombreCorto/.store in = \fAdNombreCorto,
    edad/.store in = \fAdEdad,
    icono/.store in = \fAdIcono,
    foto/.store in = \fAdFoto,
    badgeRol/.store in = \fAdBadgeRol,
    badgeTarea/.store in = \fAdBadgeTarea,
    centro/.store in = \fAdCentro,
    cita/.store in = \fAdCita,
    demographic/.store in = \fAdDemographic,
    occupation/.store in = \fAdOccupation,
    education/.store in = \fAdEducation,
    computerExp/.store in = \fAdComputerExp,
    productExp/.store in = \fAdProductExp,
    tasks/.store in = \fAdTasks,
    domainKnowledge/.store in = \fAdDomainKnowledge,
    technology/.store in = \fAdTechnology,
    attitudes/.store in = \fAdAttitudes,
    learningStyle/.store in = \fAdLearningStyle,
    criticality/.store in = \fAdCriticality,
    escenario/.store in = \fAdEscenario,
    badgesAjustes/.store in = \fAdBadgesAjustes,
}

% COMANDO PRINCIPAL \fichaAdulto
\newcommand{\fichaAdulto}[1][]{%
    \begingroup
    \pgfkeys{/fichaAdulto, default, #1}%
    % 1. BANNER SUPERIOR
    \begin{tcolorbox}[
        enhanced,
        colback=tealColor!90!black,
        colframe=tealColor!90!black,
        arc=1.6mm,
        left=7pt, right=7pt, top=2.5pt, bottom=2.5pt,
        boxsep=0pt, coltext=white
    ]
        \begin{minipage}[c]{0.70\textwidth}
            {\normalsize\textbf{\fAdBannerTitulo}}\\[0.2mm]
            {\tiny \fAdSubtitulo}
        \end{minipage}\hfill
        \begin{minipage}[c]{0.29\textwidth}
            \raggedleft\tiny
            \fAdEquipo
        \end{minipage}
    \end{tcolorbox}
    \vspace{0.7mm}
    % 2. TARJETA DE IDENTIDAD DEL ADULTO
    \begin{tcolorbox}[
        enhanced,
        colback=white,
        colframe=tealColor!70,
        arc=2.0mm,
        boxrule=0.8pt,
        left=6pt, right=6pt, top=2.8pt, bottom=2.8pt,
        boxsep=0pt,
        drop shadow=black!5
    ]
        \begin{minipage}[c]{0.10\textwidth}
            \centering
            \begin{tikzpicture}
                \ifx\fAdFoto\empty
                    \draw[fill=tealColor!8, draw=tealColor!30, line width=0.5pt, rounded corners=2mm] (0,0) rectangle (1.4, 1.55);
                    \node at (0.7, 0.95) {\color{tealColor}\scalebox{1.5}{\fAdIcono}};
                    \node[align=center] at (0.7, 0.28) {\tiny\textbf{\color{tealColor}\fAdNombreCorto}\\\fontsize{4pt}{5pt}\selectfont\color{textMuted}\fAdEdad};
                \else
                    \begin{scope}
                        \clip[rounded corners=2mm] (0,0) rectangle (1.45, 1.55);
                        \node[inner sep=0pt] at (0.725, 0.775) {\includegraphics[width=1.45cm,height=1.55cm,keepaspectratio=false]{\fAdFoto}};
                    \end{scope}
                    \draw[draw=tealColor!60, line width=0.8pt, rounded corners=2mm] (0,0) rectangle (1.45, 1.55);
                \fi
            \end{tikzpicture}
        \end{minipage}\hfill
        \begin{minipage}[c]{0.88\textwidth}
            \begin{minipage}[b]{0.68\textwidth}\raggedright
                {\large\textbf{\color{tealColor!90!black}\fAdNombre}} \hspace{2mm}
                \badge{tealColor}{\fAdBadgeRol} \badge{secondaryColor}{\fAdBadgeTarea}
            \end{minipage}\hfill
            \begin{minipage}[b]{0.31\textwidth}
                \raggedleft\small\textbf{\color{tealColor!90!black}\faSchool\ \fAdCentro}
            \end{minipage}\\[0.6mm]
            \tikz[baseline] \node[fill=tealColor!5, draw=tealColor!25, rounded corners=1.2mm, inner sep=2.5pt, text width=0.98\linewidth] {
                \small\itshape\color{tealColor!90!black}
                \faQuoteLeft\ ``\fAdCita''\ \faQuoteRight
            };
        \end{minipage}
    \end{tcolorbox}
    \vspace{0.7mm}
    % 3. LAS 11 CAJAS DE ADULTOS (5 FILAS X 2 COLS + 1 FILA COMPLETA)
    \noindent\begin{minipage}[t]{0.485\textwidth}
        \begin{tcolorbox}[cardAdultoBox, colframe=primaryColor!50, coltitle=primaryColor, title={\faIdCard\ 1. Demographic characteristics}]
            \fontsize{8.0pt}{9.4pt}\selectfont
            \begin{itemize}
                \fAdDemographic
            \end{itemize}
        \end{tcolorbox}
    \end{minipage}\hfill
    \begin{minipage}[t]{0.485\textwidth}
        \begin{tcolorbox}[cardAdultoBox, colframe=primaryColor!50, coltitle=primaryColor, title={\faUserTie\ 2. Occupation experience}]
            \fontsize{8.0pt}{9.4pt}\selectfont
            \begin{itemize}
                \fAdOccupation
            \end{itemize}
        \end{tcolorbox}
    \end{minipage}
    \vspace{0.7mm}
    \noindent\begin{minipage}[t]{0.485\textwidth}
        \begin{tcolorbox}[cardAdultoBox, colframe=tealColor!60, coltitle=tealColor!90!black, title={\faGraduationCap\ 3. Education}]
            \fontsize{8.0pt}{9.4pt}\selectfont
            \begin{itemize}
                \fAdEducation
            \end{itemize}
        \end{tcolorbox}
    \end{minipage}\hfill
    \begin{minipage}[t]{0.485\textwidth}
        \begin{tcolorbox}[cardAdultoBox, colframe=secondaryColor!60, coltitle=secondaryColor!90!black, title={\faLaptop\ 4. Computer experience}]
            \fontsize{8.0pt}{9.4pt}\selectfont
            \begin{itemize}
                \fAdComputerExp
            \end{itemize}
        \end{tcolorbox}
    \end{minipage}
    \vspace{0.7mm}
    \noindent\begin{minipage}[t]{0.485\textwidth}
        \begin{tcolorbox}[cardAdultoBox, colframe=secondaryColor!60, coltitle=secondaryColor!90!black, title={\faSearchPlus\ 5. Specific product experience}]
            \fontsize{8.0pt}{9.4pt}\selectfont
            \begin{itemize}
                \fAdProductExp
            \end{itemize}
        \end{tcolorbox}
    \end{minipage}\hfill
    \begin{minipage}[t]{0.485\textwidth}
        \begin{tcolorbox}[cardAdultoBox, colframe=greenColor!65, coltitle=greenColor!90!black, title={\faTasks\ 6. Tasks (Primary and secondary)}]
            \fontsize{8.0pt}{9.4pt}\selectfont
            \begin{itemize}
                \fAdTasks
            \end{itemize}
        \end{tcolorbox}
    \end{minipage}
    \vspace{0.7mm}
    \noindent\begin{minipage}[t]{0.485\textwidth}
        \begin{tcolorbox}[cardAdultoBox, colframe=accentColor!65, coltitle=accentColor!90!black, title={\faBrain\ 7. Domain knowledge}]
            \fontsize{8.0pt}{9.4pt}\selectfont
            \begin{itemize}
                \fAdDomainKnowledge
            \end{itemize}
        \end{tcolorbox}
    \end{minipage}\hfill
    \begin{minipage}[t]{0.485\textwidth}
        \begin{tcolorbox}[cardAdultoBox, colframe=secondaryColor!60, coltitle=secondaryColor!90!black, title={\faTools\ 8. Technology available}]
            \fontsize{8.0pt}{9.4pt}\selectfont
            \begin{itemize}
                \fAdTechnology
            \end{itemize}
        \end{tcolorbox}
    \end{minipage}
    \vspace{0.7mm}
    \noindent\begin{minipage}[t]{0.485\textwidth}
        \begin{tcolorbox}[cardAdultoBox, colframe=purpleColor!60, coltitle=purpleColor!90!black, title={\faHeart\ 9. Attitudes and values}]
            \fontsize{8.0pt}{9.4pt}\selectfont
            \begin{itemize}
                \fAdAttitudes
            \end{itemize}
        \end{tcolorbox}
    \end{minipage}\hfill
    \begin{minipage}[t]{0.485\textwidth}
        \begin{tcolorbox}[cardAdultoBox, colframe=purpleColor!60, coltitle=purpleColor!90!black, title={\faEye\ 10. Learning style}]
            \fontsize{8.0pt}{9.4pt}\selectfont
            \begin{itemize}
                \fAdLearningStyle
            \end{itemize}
        \end{tcolorbox}
    \end{minipage}
    \vspace{0.7mm}
    % FILA 6: Punto 11 (ANCHO COMPLETO — CRITICALITY OF ERRORS)
    \noindent\begin{minipage}[t]{\textwidth}
        \begin{tcolorbox}[cardAdultoWideBox, colframe=dangerColor!65, coltitle=dangerColor!90!black, title={\faExclamationTriangle\ 11. Criticality of errors (Consecuencias de un error y tolerancia del software)}]
            \fontsize{8.0pt}{9.4pt}\selectfont
            \begin{itemize}
                \fAdCriticality
            \end{itemize}
        \end{tcolorbox}
    \end{minipage}
    \vspace{0.7mm}
    % 4. BARRA INFERIOR DE INTEGRACIÓN CeRA
    \begin{tcolorbox}[
        enhanced,
        colback=tealColor!4,
        colframe=tealColor!45,
        arc=1.6mm,
        left=6pt, right=6pt, top=2.5pt, bottom=2.5pt,
        boxsep=0pt
    ]
        \scriptsize
        \textbf{\color{tealColor!90!black}\faMapSigns\ Escenario de Uso Contextual (Día típico):} \fAdEscenario\\
        \vspace{0.3mm}
        \hrule height 0.35pt \relax
        \vspace{0.5mm}
        \centerline{%
            \textbf{\color{tealColor!90!black}\faSlidersH\ Rol y Funciones en Plataforma (CeRA):} \hspace{2mm}
            \fAdBadgesAjustes%
        }
    \end{tcolorbox}
    \newpage
    \endgroup
}

% COMANDO PARA PLANTILLA EN BLANCO DE ADULTO
\newcommand{\fichaAdultoBlanco}{%
    \fichaAdulto[%
        bannerTitulo={\faUserEdit\ FICHA DE PERSONA: TUTOR/FAMILIA \textbar\ PLANTILLA DE TRABAJO},
        cita={[Cita textual representativa que exprese sus objetivos, inquietudes de gestión o necesidades respecto a la app CeRA].}%
    ]%
}

\ifdefined\isSubFile
    \expandafter\endinput
\fi

\begin{document}
\fichaAdultoBlanco
\end{document}
 y luego invocar \fichaAdulto[...] o \fichaAdultoBlanco
% ==============================================================================

\ifdefined\isSubFile
    % Si se incluye desde otro documento, los paquetes ya están cargados
\else
    \documentclass[10pt]{article}
    \usepackage[utf8]{inputenc}
    \usepackage[spanish,es-tabla]{babel}
    \usepackage[margin=0.8cm,top=0.6cm,bottom=0.6cm]{geometry}
    \usepackage{helvet}
    \renewcommand{\familydefault}{\sfdefault}
    \usepackage{microtype}
    \usepackage{graphicx}
    \usepackage{xcolor}
    \usepackage{tikz}
    \usepackage{enumitem}
    \usepackage{fontawesome5}
    \usepackage{tcolorbox}
    \tcbuselibrary{skins}
    \pagestyle{empty}
\fi

% Colores del sistema
\providecolor{primaryColor}{RGB}{20, 65, 105}     % Azul corporativo
\providecolor{secondaryColor}{RGB}{35, 115, 170}  % Azul medio accesible
\providecolor{tealColor}{RGB}{20, 130, 140}       % Verde azulado
\providecolor{greenColor}{RGB}{30, 135, 75}       % Verde bosque
\providecolor{accentColor}{RGB}{215, 105, 25}     % Naranja cálido
\providecolor{purpleColor}{RGB}{120, 60, 150}     % Púrpura suave
\providecolor{dangerColor}{RGB}{185, 45, 35}      % Rojo criticidad
\providecolor{bgCard}{RGB}{252, 253, 255}         % Fondo limpio
\providecolor{borderCard}{RGB}{215, 225, 235}     % Borde sutil
\providecolor{textDark}{RGB}{30, 40, 50}          % Texto principal oscuro
\providecolor{textMuted}{RGB}{95, 110, 125}       % Texto secundario

% Configuración de listas
\setlist[itemize]{leftmargin=3.0mm, itemsep=0.4pt, parsep=0pt, topsep=0.4pt, partopsep=0pt}

% Estilos de caja fija para Adulto
\tcbset{
    cardAdultoBox/.style={
        enhanced,
        boxrule=0.65pt,
        arc=1.8mm,
        height=27.2mm,
        left=5.0pt, right=5.0pt, top=1.5pt, bottom=1.5pt,
        boxsep=0pt,
        colback=bgCard,
        colframe=borderCard,
        fonttitle=\bfseries\small,
        fontupper=\fontsize{8.2pt}{9.6pt}\selectfont,
    },
    cardAdultoWideBox/.style={
        enhanced,
        boxrule=0.65pt,
        arc=1.8mm,
        height=20.0mm,
        left=6.0pt, right=6.0pt, top=2.0pt, bottom=2.0pt,
        boxsep=0pt,
        colback=bgCard,
        colframe=borderCard,
        fonttitle=\bfseries\small,
    }
}

% Macro para badges visuales (si no existe)
\providecommand{\badge}[2]{%
    \tikz[baseline=(b.base)] \node[fill=#1!12, text=#1!90!black, draw=#1!35, line width=0.4pt, rounded corners=1.2mm, inner sep=2.0pt, font=\bfseries\scriptsize] (b) {#2};%
}

% Definición de claves con pgfkeys para \fichaAdulto
\pgfkeys{
    /fichaAdulto/.is family, /fichaAdulto,
    default/.style = {
        bannerTitulo = {\faChalkboardTeacher\ FICHA DE PERSONA: TUTORES Y FAMILIAS (METODOLOGÍA DCU)},
        subtitulo = {Proyecto AUTOVIDA (26-175) \textbar\ Reto CeRA — Agenda Accesible PTVAL \textbar\ Asignatura MDA (UGR)},
        equipo = {\textbf{Equipo CeRA-MDA} (5 integrantes)\\\ Course 2026/2027 \textbar\ \textbf{Práctica 1}},
        nombre = {[Nombre y Apellidos del Usuario]},
        nombreCorto = {[ADULTO]},
        edad = {[Edad]},
        icono = {\faChalkboardTeacher},
        foto = {},
        badgeRol = {[Rol: Tutor / Familiar / Admin]},
        badgeTarea = {[Ámbito de Actuación]},
        centro = {CEE Purísima Concepción},
        cita = {[Cita textual representativa que exprese sus objetivos, inquietudes de gestión o necesidades respecto a la app CeRA].},
        demographic = {
            \item \textbf{Edad y Género:} [[Edad] años, [Género]].
            \item \textbf{Ubicación:} [[Ciudad / Zona residencial]].
        },
        occupation = {
            \item \textbf{Puesto actual:} [[Puesto docente, profesión o rol familiar]].
            \item \textbf{Antigüedad:} [[Años en el centro/entidad] \textbar\ [Años en puesto]].
            \item \textbf{Responsabilidades:} [[Supervisión, docencia, apoyo en el hogar]].
        },
        education = {
            \item \textbf{Titulación (Degree):} [[Nivel formativo o titulación alcanzada]].
            \item \textbf{Especialidad (Major):} [[Educación Especial, Pedagogía, etc.]].
            \item \textbf{Cursos realizados:} [[Formación en NEAE, SAAC, TIC o gestión]].
        },
        computerExp = {
            \item \textbf{Habilidades informáticas:} [[Nivel de competencia digital: medio/avanzado]].
            \item \textbf{Años de experiencia:} [[Años de uso regular de dispositivos]].
            \item \textbf{Hábitos de uso:} [[Navegación web, apps de gestión, smartphones]].
        },
        productExp = {
            \item \textbf{Productos similares:} [[Otras herramientas usadas, papel, agendas]].
            \item \textbf{Limitaciones actuales:} [[Sobrecarga de trabajo, falta de personalización]].
            \item \textbf{Expectativas:} [[Ahorro de tiempo y canal de dudas]].
        },
        tasks = {
            \item \textbf{Tareas Primarias:} 1) [[Configurar tareas]] $\rightarrow$ 2) [[Asignar a alumnos]] $\rightarrow$ 3) [[Supervisar ejecución]] $\rightarrow$ 4) [[Gestionar comedor/avisos]].
            \item \textbf{Tareas Secundarias:} [[Revisar estadísticas / chat con familias]].
        },
        domainKnowledge = {
            \item \textbf{Comprensión del área:} [[Conocimiento sobre necesidades NEAE y rutinas]].
            \item \textbf{Conocimiento del usuario:} [[Grado de personalización que requiere cada estudiante]].
        },
        technology = {
            \item \textbf{Hardware:} [[Ordenador portátil / tablet / móvil Android]].
            \item \textbf{Software:} [[Navegador web / app móvil]].
            \item \textbf{Herramientas:} [[Conexión a internet]].
        },
        attitudes = {
            \item \textbf{Preferencias:} [[Sencillez, rapidez operativa y fiabilidad en tiempo real]].
            \item \textbf{Miedo a la tecnología:} [[Preocupación por pérdida de datos o fallos de red]].
            \item \textbf{Valores:} [[Foco en el progreso y bienestar del alumnado]].
        },
        learningStyle = {
            \item \textbf{Perfil de aprendizaje:} [[\textbf{Práctico} / \textbf{Estructurado} / \textbf{Visual}]].
            \item \textbf{Canales óptimos:} [[Tutoriales cortos, interfaz intuitiva sin manual]].
            \item \textbf{Curva de adopción:} [[Inmediata si soluciona problemas reales]].
        },
        criticality = {
            \item \textbf{Consecuencias de un error:} \textit{Logístico:} [[¿Dietas erróneas, tareas desincronizadas?]] \textit{Pedagógico:} [[¿Confusión en el horario del alumno, estrés?]]
            \item \textbf{Tolerancia del software:} [[Alertas claras ante cambios sensibles, confirmación de guardado]].
        },
        escenario = {[[Describir brevemente en 3-4 líneas cómo accede al panel de la app, planifica o revisa las tareas de sus estudiantes y cómo la herramienta le facilita la jornada eliminando fricciones organizativas.]]},
        badgesAjustes = {
            \badge{tealColor}{\faDesktop\ [[Panel Web]]} \hspace{2mm}
            \badge{primaryColor}{\faCalendarCheck\ [[Planificación]]} \hspace{2mm}
            \badge{accentColor}{\faUtensils\ [[Comedor/Dietas]]} \hspace{2mm}
            \badge{greenColor}{\faCheckCircle\ [[Supervisión]]}%
        }
    },
    bannerTitulo/.store in = \fAdBannerTitulo,
    subtitulo/.store in = \fAdSubtitulo,
    equipo/.store in = \fAdEquipo,
    nombre/.store in = \fAdNombre,
    nombreCorto/.store in = \fAdNombreCorto,
    edad/.store in = \fAdEdad,
    icono/.store in = \fAdIcono,
    foto/.store in = \fAdFoto,
    badgeRol/.store in = \fAdBadgeRol,
    badgeTarea/.store in = \fAdBadgeTarea,
    centro/.store in = \fAdCentro,
    cita/.store in = \fAdCita,
    demographic/.store in = \fAdDemographic,
    occupation/.store in = \fAdOccupation,
    education/.store in = \fAdEducation,
    computerExp/.store in = \fAdComputerExp,
    productExp/.store in = \fAdProductExp,
    tasks/.store in = \fAdTasks,
    domainKnowledge/.store in = \fAdDomainKnowledge,
    technology/.store in = \fAdTechnology,
    attitudes/.store in = \fAdAttitudes,
    learningStyle/.store in = \fAdLearningStyle,
    criticality/.store in = \fAdCriticality,
    escenario/.store in = \fAdEscenario,
    badgesAjustes/.store in = \fAdBadgesAjustes,
}

% COMANDO PRINCIPAL \fichaAdulto
\newcommand{\fichaAdulto}[1][]{%
    \begingroup
    \pgfkeys{/fichaAdulto, default, #1}%
    % 1. BANNER SUPERIOR
    \begin{tcolorbox}[
        enhanced,
        colback=tealColor!90!black,
        colframe=tealColor!90!black,
        arc=1.6mm,
        left=7pt, right=7pt, top=2.5pt, bottom=2.5pt,
        boxsep=0pt, coltext=white
    ]
        \begin{minipage}[c]{0.70\textwidth}
            {\normalsize\textbf{\fAdBannerTitulo}}\\[0.2mm]
            {\tiny \fAdSubtitulo}
        \end{minipage}\hfill
        \begin{minipage}[c]{0.29\textwidth}
            \raggedleft\tiny
            \fAdEquipo
        \end{minipage}
    \end{tcolorbox}
    \vspace{0.7mm}
    % 2. TARJETA DE IDENTIDAD DEL ADULTO
    \begin{tcolorbox}[
        enhanced,
        colback=white,
        colframe=tealColor!70,
        arc=2.0mm,
        boxrule=0.8pt,
        left=6pt, right=6pt, top=2.8pt, bottom=2.8pt,
        boxsep=0pt,
        drop shadow=black!5
    ]
        \begin{minipage}[c]{0.10\textwidth}
            \centering
            \begin{tikzpicture}
                \ifx\fAdFoto\empty
                    \draw[fill=tealColor!8, draw=tealColor!30, line width=0.5pt, rounded corners=2mm] (0,0) rectangle (1.4, 1.55);
                    \node at (0.7, 0.95) {\color{tealColor}\scalebox{1.5}{\fAdIcono}};
                    \node[align=center] at (0.7, 0.28) {\tiny\textbf{\color{tealColor}\fAdNombreCorto}\\\fontsize{4pt}{5pt}\selectfont\color{textMuted}\fAdEdad};
                \else
                    \begin{scope}
                        \clip[rounded corners=2mm] (0,0) rectangle (1.45, 1.55);
                        \node[inner sep=0pt] at (0.725, 0.775) {\includegraphics[width=1.45cm,height=1.55cm,keepaspectratio=false]{\fAdFoto}};
                    \end{scope}
                    \draw[draw=tealColor!60, line width=0.8pt, rounded corners=2mm] (0,0) rectangle (1.45, 1.55);
                \fi
            \end{tikzpicture}
        \end{minipage}\hfill
        \begin{minipage}[c]{0.88\textwidth}
            \begin{minipage}[b]{0.68\textwidth}\raggedright
                {\large\textbf{\color{tealColor!90!black}\fAdNombre}} \hspace{2mm}
                \badge{tealColor}{\fAdBadgeRol} \badge{secondaryColor}{\fAdBadgeTarea}
            \end{minipage}\hfill
            \begin{minipage}[b]{0.31\textwidth}
                \raggedleft\small\textbf{\color{tealColor!90!black}\faSchool\ \fAdCentro}
            \end{minipage}\\[0.6mm]
            \tikz[baseline] \node[fill=tealColor!5, draw=tealColor!25, rounded corners=1.2mm, inner sep=2.5pt, text width=0.98\linewidth] {
                \small\itshape\color{tealColor!90!black}
                \faQuoteLeft\ ``\fAdCita''\ \faQuoteRight
            };
        \end{minipage}
    \end{tcolorbox}
    \vspace{0.7mm}
    % 3. LAS 11 CAJAS DE ADULTOS (5 FILAS X 2 COLS + 1 FILA COMPLETA)
    \noindent\begin{minipage}[t]{0.485\textwidth}
        \begin{tcolorbox}[cardAdultoBox, colframe=primaryColor!50, coltitle=primaryColor, title={\faIdCard\ 1. Demographic characteristics}]
            \fontsize{8.0pt}{9.4pt}\selectfont
            \begin{itemize}
                \fAdDemographic
            \end{itemize}
        \end{tcolorbox}
    \end{minipage}\hfill
    \begin{minipage}[t]{0.485\textwidth}
        \begin{tcolorbox}[cardAdultoBox, colframe=primaryColor!50, coltitle=primaryColor, title={\faUserTie\ 2. Occupation experience}]
            \fontsize{8.0pt}{9.4pt}\selectfont
            \begin{itemize}
                \fAdOccupation
            \end{itemize}
        \end{tcolorbox}
    \end{minipage}
    \vspace{0.7mm}
    \noindent\begin{minipage}[t]{0.485\textwidth}
        \begin{tcolorbox}[cardAdultoBox, colframe=tealColor!60, coltitle=tealColor!90!black, title={\faGraduationCap\ 3. Education}]
            \fontsize{8.0pt}{9.4pt}\selectfont
            \begin{itemize}
                \fAdEducation
            \end{itemize}
        \end{tcolorbox}
    \end{minipage}\hfill
    \begin{minipage}[t]{0.485\textwidth}
        \begin{tcolorbox}[cardAdultoBox, colframe=secondaryColor!60, coltitle=secondaryColor!90!black, title={\faLaptop\ 4. Computer experience}]
            \fontsize{8.0pt}{9.4pt}\selectfont
            \begin{itemize}
                \fAdComputerExp
            \end{itemize}
        \end{tcolorbox}
    \end{minipage}
    \vspace{0.7mm}
    \noindent\begin{minipage}[t]{0.485\textwidth}
        \begin{tcolorbox}[cardAdultoBox, colframe=secondaryColor!60, coltitle=secondaryColor!90!black, title={\faSearchPlus\ 5. Specific product experience}]
            \fontsize{8.0pt}{9.4pt}\selectfont
            \begin{itemize}
                \fAdProductExp
            \end{itemize}
        \end{tcolorbox}
    \end{minipage}\hfill
    \begin{minipage}[t]{0.485\textwidth}
        \begin{tcolorbox}[cardAdultoBox, colframe=greenColor!65, coltitle=greenColor!90!black, title={\faTasks\ 6. Tasks (Primary and secondary)}]
            \fontsize{8.0pt}{9.4pt}\selectfont
            \begin{itemize}
                \fAdTasks
            \end{itemize}
        \end{tcolorbox}
    \end{minipage}
    \vspace{0.7mm}
    \noindent\begin{minipage}[t]{0.485\textwidth}
        \begin{tcolorbox}[cardAdultoBox, colframe=accentColor!65, coltitle=accentColor!90!black, title={\faBrain\ 7. Domain knowledge}]
            \fontsize{8.0pt}{9.4pt}\selectfont
            \begin{itemize}
                \fAdDomainKnowledge
            \end{itemize}
        \end{tcolorbox}
    \end{minipage}\hfill
    \begin{minipage}[t]{0.485\textwidth}
        \begin{tcolorbox}[cardAdultoBox, colframe=secondaryColor!60, coltitle=secondaryColor!90!black, title={\faTools\ 8. Technology available}]
            \fontsize{8.0pt}{9.4pt}\selectfont
            \begin{itemize}
                \fAdTechnology
            \end{itemize}
        \end{tcolorbox}
    \end{minipage}
    \vspace{0.7mm}
    \noindent\begin{minipage}[t]{0.485\textwidth}
        \begin{tcolorbox}[cardAdultoBox, colframe=purpleColor!60, coltitle=purpleColor!90!black, title={\faHeart\ 9. Attitudes and values}]
            \fontsize{8.0pt}{9.4pt}\selectfont
            \begin{itemize}
                \fAdAttitudes
            \end{itemize}
        \end{tcolorbox}
    \end{minipage}\hfill
    \begin{minipage}[t]{0.485\textwidth}
        \begin{tcolorbox}[cardAdultoBox, colframe=purpleColor!60, coltitle=purpleColor!90!black, title={\faEye\ 10. Learning style}]
            \fontsize{8.0pt}{9.4pt}\selectfont
            \begin{itemize}
                \fAdLearningStyle
            \end{itemize}
        \end{tcolorbox}
    \end{minipage}
    \vspace{0.7mm}
    % FILA 6: Punto 11 (ANCHO COMPLETO — CRITICALITY OF ERRORS)
    \noindent\begin{minipage}[t]{\textwidth}
        \begin{tcolorbox}[cardAdultoWideBox, colframe=dangerColor!65, coltitle=dangerColor!90!black, title={\faExclamationTriangle\ 11. Criticality of errors (Consecuencias de un error y tolerancia del software)}]
            \fontsize{8.0pt}{9.4pt}\selectfont
            \begin{itemize}
                \fAdCriticality
            \end{itemize}
        \end{tcolorbox}
    \end{minipage}
    \vspace{0.7mm}
    % 4. BARRA INFERIOR DE INTEGRACIÓN CeRA
    \begin{tcolorbox}[
        enhanced,
        colback=tealColor!4,
        colframe=tealColor!45,
        arc=1.6mm,
        left=6pt, right=6pt, top=2.5pt, bottom=2.5pt,
        boxsep=0pt
    ]
        \scriptsize
        \textbf{\color{tealColor!90!black}\faMapSigns\ Escenario de Uso Contextual (Día típico):} \fAdEscenario\\
        \vspace{0.3mm}
        \hrule height 0.35pt \relax
        \vspace{0.5mm}
        \centerline{%
            \textbf{\color{tealColor!90!black}\faSlidersH\ Rol y Funciones en Plataforma (CeRA):} \hspace{2mm}
            \fAdBadgesAjustes%
        }
    \end{tcolorbox}
    \newpage
    \endgroup
}

% COMANDO PARA PLANTILLA EN BLANCO DE ADULTO
\newcommand{\fichaAdultoBlanco}{%
    \fichaAdulto[%
        bannerTitulo={\faUserEdit\ FICHA DE PERSONA: TUTOR/FAMILIA \textbar\ PLANTILLA DE TRABAJO},
        cita={[Cita textual representativa que exprese sus objetivos, inquietudes de gestión o necesidades respecto a la app CeRA].}%
    ]%
}

\ifdefined\isSubFile
    \expandafter\endinput
\fi

\begin{document}
\fichaAdultoBlanco
\end{document}
 y luego invocar \fichaAdulto[...] o \fichaAdultoBlanco
% ==============================================================================

\ifdefined\isSubFile
    % Si se incluye desde otro documento, los paquetes ya están cargados
\else
    \documentclass[10pt]{article}
    \usepackage[utf8]{inputenc}
    \usepackage[spanish,es-tabla]{babel}
    \usepackage[margin=0.8cm,top=0.6cm,bottom=0.6cm]{geometry}
    \usepackage{helvet}
    \renewcommand{\familydefault}{\sfdefault}
    \usepackage{microtype}
    \usepackage{graphicx}
    \usepackage{xcolor}
    \usepackage{tikz}
    \usepackage{enumitem}
    \usepackage{fontawesome5}
    \usepackage{tcolorbox}
    \tcbuselibrary{skins}
    \pagestyle{empty}
\fi

% Colores del sistema
\providecolor{primaryColor}{RGB}{20, 65, 105}     % Azul corporativo
\providecolor{secondaryColor}{RGB}{35, 115, 170}  % Azul medio accesible
\providecolor{tealColor}{RGB}{20, 130, 140}       % Verde azulado
\providecolor{greenColor}{RGB}{30, 135, 75}       % Verde bosque
\providecolor{accentColor}{RGB}{215, 105, 25}     % Naranja cálido
\providecolor{purpleColor}{RGB}{120, 60, 150}     % Púrpura suave
\providecolor{dangerColor}{RGB}{185, 45, 35}      % Rojo criticidad
\providecolor{bgCard}{RGB}{252, 253, 255}         % Fondo limpio
\providecolor{borderCard}{RGB}{215, 225, 235}     % Borde sutil
\providecolor{textDark}{RGB}{30, 40, 50}          % Texto principal oscuro
\providecolor{textMuted}{RGB}{95, 110, 125}       % Texto secundario

% Configuración de listas
\setlist[itemize]{leftmargin=3.0mm, itemsep=0.4pt, parsep=0pt, topsep=0.4pt, partopsep=0pt}

% Estilos de caja fija para Adulto
\tcbset{
    cardAdultoBox/.style={
        enhanced,
        boxrule=0.65pt,
        arc=1.8mm,
        height=27.2mm,
        left=5.0pt, right=5.0pt, top=1.5pt, bottom=1.5pt,
        boxsep=0pt,
        colback=bgCard,
        colframe=borderCard,
        fonttitle=\bfseries\small,
        fontupper=\fontsize{8.2pt}{9.6pt}\selectfont,
    },
    cardAdultoWideBox/.style={
        enhanced,
        boxrule=0.65pt,
        arc=1.8mm,
        height=20.0mm,
        left=6.0pt, right=6.0pt, top=2.0pt, bottom=2.0pt,
        boxsep=0pt,
        colback=bgCard,
        colframe=borderCard,
        fonttitle=\bfseries\small,
    }
}

% Macro para badges visuales (si no existe)
\providecommand{\badge}[2]{%
    \tikz[baseline=(b.base)] \node[fill=#1!12, text=#1!90!black, draw=#1!35, line width=0.4pt, rounded corners=1.2mm, inner sep=2.0pt, font=\bfseries\scriptsize] (b) {#2};%
}

% Definición de claves con pgfkeys para \fichaAdulto
\pgfkeys{
    /fichaAdulto/.is family, /fichaAdulto,
    default/.style = {
        bannerTitulo = {\faChalkboardTeacher\ FICHA DE PERSONA: TUTORES Y FAMILIAS (METODOLOGÍA DCU)},
        subtitulo = {Proyecto AUTOVIDA (26-175) \textbar\ Reto CeRA — Agenda Accesible PTVAL \textbar\ Asignatura MDA (UGR)},
        equipo = {\textbf{Equipo CeRA-MDA} (5 integrantes)\\\ Course 2026/2027 \textbar\ \textbf{Práctica 1}},
        nombre = {[Nombre y Apellidos del Usuario]},
        nombreCorto = {[ADULTO]},
        edad = {[Edad]},
        icono = {\faChalkboardTeacher},
        foto = {},
        badgeRol = {[Rol: Tutor / Familiar / Admin]},
        badgeTarea = {[Ámbito de Actuación]},
        centro = {CEE Purísima Concepción},
        cita = {[Cita textual representativa que exprese sus objetivos, inquietudes de gestión o necesidades respecto a la app CeRA].},
        demographic = {
            \item \textbf{Edad y Género:} [[Edad] años, [Género]].
            \item \textbf{Ubicación:} [[Ciudad / Zona residencial]].
        },
        occupation = {
            \item \textbf{Puesto actual:} [[Puesto docente, profesión o rol familiar]].
            \item \textbf{Antigüedad:} [[Años en el centro/entidad] \textbar\ [Años en puesto]].
            \item \textbf{Responsabilidades:} [[Supervisión, docencia, apoyo en el hogar]].
        },
        education = {
            \item \textbf{Titulación (Degree):} [[Nivel formativo o titulación alcanzada]].
            \item \textbf{Especialidad (Major):} [[Educación Especial, Pedagogía, etc.]].
            \item \textbf{Cursos realizados:} [[Formación en NEAE, SAAC, TIC o gestión]].
        },
        computerExp = {
            \item \textbf{Habilidades informáticas:} [[Nivel de competencia digital: medio/avanzado]].
            \item \textbf{Años de experiencia:} [[Años de uso regular de dispositivos]].
            \item \textbf{Hábitos de uso:} [[Navegación web, apps de gestión, smartphones]].
        },
        productExp = {
            \item \textbf{Productos similares:} [[Otras herramientas usadas, papel, agendas]].
            \item \textbf{Limitaciones actuales:} [[Sobrecarga de trabajo, falta de personalización]].
            \item \textbf{Expectativas:} [[Ahorro de tiempo y canal de dudas]].
        },
        tasks = {
            \item \textbf{Tareas Primarias:} 1) [[Configurar tareas]] $\rightarrow$ 2) [[Asignar a alumnos]] $\rightarrow$ 3) [[Supervisar ejecución]] $\rightarrow$ 4) [[Gestionar comedor/avisos]].
            \item \textbf{Tareas Secundarias:} [[Revisar estadísticas / chat con familias]].
        },
        domainKnowledge = {
            \item \textbf{Comprensión del área:} [[Conocimiento sobre necesidades NEAE y rutinas]].
            \item \textbf{Conocimiento del usuario:} [[Grado de personalización que requiere cada estudiante]].
        },
        technology = {
            \item \textbf{Hardware:} [[Ordenador portátil / tablet / móvil Android]].
            \item \textbf{Software:} [[Navegador web / app móvil]].
            \item \textbf{Herramientas:} [[Conexión a internet]].
        },
        attitudes = {
            \item \textbf{Preferencias:} [[Sencillez, rapidez operativa y fiabilidad en tiempo real]].
            \item \textbf{Miedo a la tecnología:} [[Preocupación por pérdida de datos o fallos de red]].
            \item \textbf{Valores:} [[Foco en el progreso y bienestar del alumnado]].
        },
        learningStyle = {
            \item \textbf{Perfil de aprendizaje:} [[\textbf{Práctico} / \textbf{Estructurado} / \textbf{Visual}]].
            \item \textbf{Canales óptimos:} [[Tutoriales cortos, interfaz intuitiva sin manual]].
            \item \textbf{Curva de adopción:} [[Inmediata si soluciona problemas reales]].
        },
        criticality = {
            \item \textbf{Consecuencias de un error:} \textit{Logístico:} [[¿Dietas erróneas, tareas desincronizadas?]] \textit{Pedagógico:} [[¿Confusión en el horario del alumno, estrés?]]
            \item \textbf{Tolerancia del software:} [[Alertas claras ante cambios sensibles, confirmación de guardado]].
        },
        escenario = {[[Describir brevemente en 3-4 líneas cómo accede al panel de la app, planifica o revisa las tareas de sus estudiantes y cómo la herramienta le facilita la jornada eliminando fricciones organizativas.]]},
        badgesAjustes = {
            \badge{tealColor}{\faDesktop\ [[Panel Web]]} \hspace{2mm}
            \badge{primaryColor}{\faCalendarCheck\ [[Planificación]]} \hspace{2mm}
            \badge{accentColor}{\faUtensils\ [[Comedor/Dietas]]} \hspace{2mm}
            \badge{greenColor}{\faCheckCircle\ [[Supervisión]]}%
        }
    },
    bannerTitulo/.store in = \fAdBannerTitulo,
    subtitulo/.store in = \fAdSubtitulo,
    equipo/.store in = \fAdEquipo,
    nombre/.store in = \fAdNombre,
    nombreCorto/.store in = \fAdNombreCorto,
    edad/.store in = \fAdEdad,
    icono/.store in = \fAdIcono,
    foto/.store in = \fAdFoto,
    badgeRol/.store in = \fAdBadgeRol,
    badgeTarea/.store in = \fAdBadgeTarea,
    centro/.store in = \fAdCentro,
    cita/.store in = \fAdCita,
    demographic/.store in = \fAdDemographic,
    occupation/.store in = \fAdOccupation,
    education/.store in = \fAdEducation,
    computerExp/.store in = \fAdComputerExp,
    productExp/.store in = \fAdProductExp,
    tasks/.store in = \fAdTasks,
    domainKnowledge/.store in = \fAdDomainKnowledge,
    technology/.store in = \fAdTechnology,
    attitudes/.store in = \fAdAttitudes,
    learningStyle/.store in = \fAdLearningStyle,
    criticality/.store in = \fAdCriticality,
    escenario/.store in = \fAdEscenario,
    badgesAjustes/.store in = \fAdBadgesAjustes,
}

% COMANDO PRINCIPAL \fichaAdulto
\newcommand{\fichaAdulto}[1][]{%
    \begingroup
    \pgfkeys{/fichaAdulto, default, #1}%
    % 1. BANNER SUPERIOR
    \begin{tcolorbox}[
        enhanced,
        colback=tealColor!90!black,
        colframe=tealColor!90!black,
        arc=1.6mm,
        left=7pt, right=7pt, top=2.5pt, bottom=2.5pt,
        boxsep=0pt, coltext=white
    ]
        \begin{minipage}[c]{0.70\textwidth}
            {\normalsize\textbf{\fAdBannerTitulo}}\\[0.2mm]
            {\tiny \fAdSubtitulo}
        \end{minipage}\hfill
        \begin{minipage}[c]{0.29\textwidth}
            \raggedleft\tiny
            \fAdEquipo
        \end{minipage}
    \end{tcolorbox}
    \vspace{0.7mm}
    % 2. TARJETA DE IDENTIDAD DEL ADULTO
    \begin{tcolorbox}[
        enhanced,
        colback=white,
        colframe=tealColor!70,
        arc=2.0mm,
        boxrule=0.8pt,
        left=6pt, right=6pt, top=2.8pt, bottom=2.8pt,
        boxsep=0pt,
        drop shadow=black!5
    ]
        \begin{minipage}[c]{0.10\textwidth}
            \centering
            \begin{tikzpicture}
                \ifx\fAdFoto\empty
                    \draw[fill=tealColor!8, draw=tealColor!30, line width=0.5pt, rounded corners=2mm] (0,0) rectangle (1.4, 1.55);
                    \node at (0.7, 0.95) {\color{tealColor}\scalebox{1.5}{\fAdIcono}};
                    \node[align=center] at (0.7, 0.28) {\tiny\textbf{\color{tealColor}\fAdNombreCorto}\\\fontsize{4pt}{5pt}\selectfont\color{textMuted}\fAdEdad};
                \else
                    \begin{scope}
                        \clip[rounded corners=2mm] (0,0) rectangle (1.45, 1.55);
                        \node[inner sep=0pt] at (0.725, 0.775) {\includegraphics[width=1.45cm,height=1.55cm,keepaspectratio=false]{\fAdFoto}};
                    \end{scope}
                    \draw[draw=tealColor!60, line width=0.8pt, rounded corners=2mm] (0,0) rectangle (1.45, 1.55);
                \fi
            \end{tikzpicture}
        \end{minipage}\hfill
        \begin{minipage}[c]{0.88\textwidth}
            \begin{minipage}[b]{0.68\textwidth}\raggedright
                {\large\textbf{\color{tealColor!90!black}\fAdNombre}} \hspace{2mm}
                \badge{tealColor}{\fAdBadgeRol} \badge{secondaryColor}{\fAdBadgeTarea}
            \end{minipage}\hfill
            \begin{minipage}[b]{0.31\textwidth}
                \raggedleft\small\textbf{\color{tealColor!90!black}\faSchool\ \fAdCentro}
            \end{minipage}\\[0.6mm]
            \tikz[baseline] \node[fill=tealColor!5, draw=tealColor!25, rounded corners=1.2mm, inner sep=2.5pt, text width=0.98\linewidth] {
                \small\itshape\color{tealColor!90!black}
                \faQuoteLeft\ ``\fAdCita''\ \faQuoteRight
            };
        \end{minipage}
    \end{tcolorbox}
    \vspace{0.7mm}
    % 3. LAS 11 CAJAS DE ADULTOS (5 FILAS X 2 COLS + 1 FILA COMPLETA)
    \noindent\begin{minipage}[t]{0.485\textwidth}
        \begin{tcolorbox}[cardAdultoBox, colframe=primaryColor!50, coltitle=primaryColor, title={\faIdCard\ 1. Demographic characteristics}]
            \fontsize{8.0pt}{9.4pt}\selectfont
            \begin{itemize}
                \fAdDemographic
            \end{itemize}
        \end{tcolorbox}
    \end{minipage}\hfill
    \begin{minipage}[t]{0.485\textwidth}
        \begin{tcolorbox}[cardAdultoBox, colframe=primaryColor!50, coltitle=primaryColor, title={\faUserTie\ 2. Occupation experience}]
            \fontsize{8.0pt}{9.4pt}\selectfont
            \begin{itemize}
                \fAdOccupation
            \end{itemize}
        \end{tcolorbox}
    \end{minipage}
    \vspace{0.7mm}
    \noindent\begin{minipage}[t]{0.485\textwidth}
        \begin{tcolorbox}[cardAdultoBox, colframe=tealColor!60, coltitle=tealColor!90!black, title={\faGraduationCap\ 3. Education}]
            \fontsize{8.0pt}{9.4pt}\selectfont
            \begin{itemize}
                \fAdEducation
            \end{itemize}
        \end{tcolorbox}
    \end{minipage}\hfill
    \begin{minipage}[t]{0.485\textwidth}
        \begin{tcolorbox}[cardAdultoBox, colframe=secondaryColor!60, coltitle=secondaryColor!90!black, title={\faLaptop\ 4. Computer experience}]
            \fontsize{8.0pt}{9.4pt}\selectfont
            \begin{itemize}
                \fAdComputerExp
            \end{itemize}
        \end{tcolorbox}
    \end{minipage}
    \vspace{0.7mm}
    \noindent\begin{minipage}[t]{0.485\textwidth}
        \begin{tcolorbox}[cardAdultoBox, colframe=secondaryColor!60, coltitle=secondaryColor!90!black, title={\faSearchPlus\ 5. Specific product experience}]
            \fontsize{8.0pt}{9.4pt}\selectfont
            \begin{itemize}
                \fAdProductExp
            \end{itemize}
        \end{tcolorbox}
    \end{minipage}\hfill
    \begin{minipage}[t]{0.485\textwidth}
        \begin{tcolorbox}[cardAdultoBox, colframe=greenColor!65, coltitle=greenColor!90!black, title={\faTasks\ 6. Tasks (Primary and secondary)}]
            \fontsize{8.0pt}{9.4pt}\selectfont
            \begin{itemize}
                \fAdTasks
            \end{itemize}
        \end{tcolorbox}
    \end{minipage}
    \vspace{0.7mm}
    \noindent\begin{minipage}[t]{0.485\textwidth}
        \begin{tcolorbox}[cardAdultoBox, colframe=accentColor!65, coltitle=accentColor!90!black, title={\faBrain\ 7. Domain knowledge}]
            \fontsize{8.0pt}{9.4pt}\selectfont
            \begin{itemize}
                \fAdDomainKnowledge
            \end{itemize}
        \end{tcolorbox}
    \end{minipage}\hfill
    \begin{minipage}[t]{0.485\textwidth}
        \begin{tcolorbox}[cardAdultoBox, colframe=secondaryColor!60, coltitle=secondaryColor!90!black, title={\faTools\ 8. Technology available}]
            \fontsize{8.0pt}{9.4pt}\selectfont
            \begin{itemize}
                \fAdTechnology
            \end{itemize}
        \end{tcolorbox}
    \end{minipage}
    \vspace{0.7mm}
    \noindent\begin{minipage}[t]{0.485\textwidth}
        \begin{tcolorbox}[cardAdultoBox, colframe=purpleColor!60, coltitle=purpleColor!90!black, title={\faHeart\ 9. Attitudes and values}]
            \fontsize{8.0pt}{9.4pt}\selectfont
            \begin{itemize}
                \fAdAttitudes
            \end{itemize}
        \end{tcolorbox}
    \end{minipage}\hfill
    \begin{minipage}[t]{0.485\textwidth}
        \begin{tcolorbox}[cardAdultoBox, colframe=purpleColor!60, coltitle=purpleColor!90!black, title={\faEye\ 10. Learning style}]
            \fontsize{8.0pt}{9.4pt}\selectfont
            \begin{itemize}
                \fAdLearningStyle
            \end{itemize}
        \end{tcolorbox}
    \end{minipage}
    \vspace{0.7mm}
    % FILA 6: Punto 11 (ANCHO COMPLETO — CRITICALITY OF ERRORS)
    \noindent\begin{minipage}[t]{\textwidth}
        \begin{tcolorbox}[cardAdultoWideBox, colframe=dangerColor!65, coltitle=dangerColor!90!black, title={\faExclamationTriangle\ 11. Criticality of errors (Consecuencias de un error y tolerancia del software)}]
            \fontsize{8.0pt}{9.4pt}\selectfont
            \begin{itemize}
                \fAdCriticality
            \end{itemize}
        \end{tcolorbox}
    \end{minipage}
    \vspace{0.7mm}
    % 4. BARRA INFERIOR DE INTEGRACIÓN CeRA
    \begin{tcolorbox}[
        enhanced,
        colback=tealColor!4,
        colframe=tealColor!45,
        arc=1.6mm,
        left=6pt, right=6pt, top=2.5pt, bottom=2.5pt,
        boxsep=0pt
    ]
        \scriptsize
        \textbf{\color{tealColor!90!black}\faMapSigns\ Escenario de Uso Contextual (Día típico):} \fAdEscenario\\
        \vspace{0.3mm}
        \hrule height 0.35pt \relax
        \vspace{0.5mm}
        \centerline{%
            \textbf{\color{tealColor!90!black}\faSlidersH\ Rol y Funciones en Plataforma (CeRA):} \hspace{2mm}
            \fAdBadgesAjustes%
        }
    \end{tcolorbox}
    \newpage
    \endgroup
}

% COMANDO PARA PLANTILLA EN BLANCO DE ADULTO
\newcommand{\fichaAdultoBlanco}{%
    \fichaAdulto[%
        bannerTitulo={\faUserEdit\ FICHA DE PERSONA: TUTOR/FAMILIA \textbar\ PLANTILLA DE TRABAJO},
        cita={[Cita textual representativa que exprese sus objetivos, inquietudes de gestión o necesidades respecto a la app CeRA].}%
    ]%
}

\ifdefined\isSubFile
    \expandafter\endinput
\fi

\begin{document}
\fichaAdultoBlanco
\end{document}


\begin{document}

% ==============================================================================
% PERSONA 1: ALUMNO CON DISCAPACIDAD ÚNICAMENTE MOTRIZ
% ==============================================================================
\fichaAlumno[%
    nombre={Carlos Martínez Ruiz},
    nombreCorto={CARLOS M.},
    edad={18 años},
    icono={\faWheelchair},
    foto={img/carlos_martinez.jpg},
    badgeRol={Alumno PTVAL},
    badgeTarea={Tareas Rotativas},
    centro={CEE Purísima Concepción},
    cita={Comprendo todo lo que hay que hacer y quiero ser independiente, pero mis manos no me obedecen con pantallas táctiles pequeñas; necesito botones grandes o un pulsador externo.},
    demographic={%
        \item \textbf{Edad y Género (Age / Gender):} 18 años, Varón.
        \item \textbf{Ubicación (Location):} Granada capital (domicilio familiar).
        \item \textbf{Perfil diagnóstico:} Parálisis cerebral con tetraparesia espástica. Usuario de silla de ruedas. Capacidad cognitiva, razonamiento y lenguaje normotípicos (sin discapacidad intelectual).
    },
    computerExp={%
        \item \textbf{Habilidades informáticas:} Manejo táctil asistido y uso de joystick en su silla. No realiza pinza fina ni tecleo rápido.
        \item \textbf{Años de experiencia:} 6 años manejando tablets y ordenadores adaptados.
        \item \textbf{Hábitos de uso:} Navegación web, dictado por voz, consumo de vídeos y mensajería adaptada.
        \item \textbf{Autonomía digital:} Muy alta si la app ofrece botones grandes o pulsador.
    },
    productExp={%
        \item \textbf{Productos similares:} Ha probado agendas digitales estándar en móvil.
        \item \textbf{Limitaciones actuales:} Los botones pequeños y los gestos táctiles le provocan toques involuntarios por espasmos.
        \item \textbf{Tendencias de uso:} Prefiere interfaces limpias con opción de pulsador adaptado.
    },
    tasks={%
        \item \textbf{Tareas Primarias:} 
        1) Login con botón grande $\rightarrow$ 2) Consultar la tarea rotativa del día (visitas, comedor, material) $\rightarrow$ 3) Ejecutar los pasos guiados con la tablet $\rightarrow$ 4) Confirmar tarea con pulsador.
        \item \textbf{Tareas Secundarias:} Avisar al tutor si un ascensor o puerta están bloqueados.
    },
    domainKnowledge={%
        \item \textbf{Comprensión del área:} Conoce al detalle todas las aulas, talleres y personal del centro.
        \item \textbf{Comprensión temporal:} Completa; entiende horas y minutos estándar, pero agradece un temporizador gráfico.
        \item \textbf{Destreza procedimental:} Gran capacidad para realizar tareas del centro con apoyo técnico accesible.
    },
    technology={%
        \item \textbf{Hardware:} Tablet Android y pulsador externo.
        \item \textbf{Software:} Sistema Android y lector de códigos QR.
        \item \textbf{Herramientas habituales:} Conexión Wi-Fi.
    },
    attitudes={%
        \item \textbf{Preferencias:} Interfaz clara, adulta y no infantilizada, con respuesta rápida.
        \item \textbf{Miedo a la tecnología:} Ninguno; siente frustración si la interfaz es pequeña y le obliga a pedir ayuda física.
        \item \textbf{Valores:} Orgullo por su autonomía y por participar en las tareas rotativas del centro.
    },
    learningStyle={%
        \item \textbf{Perfil de aprendizaje:} Cognitivo estándar (\textbf{Visual y Auditivo}).
        \item \textbf{Canales óptimos:} Iconos claros, texto legible, vídeos explicativos y sonido de confirmación.
        \item \textbf{Metodología:} Asimilación directa; requiere únicamente adaptación física de entrada.
    },
    criticality={%
        \item \textbf{Consecuencias de un error:} \textit{Emocional:} Frustración e impotencia al no poder pulsar el control deseado. \textit{Operativo:} Pulsación involuntaria por espasmo que marque una tarea sin terminar.
        \item \textbf{Tolerancia del software:} Filtro antitemblor (retardo de pulsación), botones grandes y confirmación antes de desmarcar tareas.
    },
    escenario={A las 11:00h, Carlos recibe un aviso en su tablet. La app le indica su tarea rotativa de hoy: acompañar al grupo de visitas escolares. Conduce su silla a la entrada, escanea el código QR de cada aula con la tablet para reproducir la información y presiona su pulsador para marcar cada hito como completado.},
    badgesAjustes={%
        \badge{primaryColor}{\faHandPointer\ Pulsador Externo} \hspace{2mm}
        \badge{secondaryColor}{\faExpand\ Botones Grandes} \hspace{2mm}
        \badge{accentColor}{\faSync\ Tareas Rotativas} \hspace{2mm}
        \badge{purpleColor}{\faLock\ Filtro Antitemblor}%
    }%
]

% ==============================================================================
% PERSONA 2: ALUMNO CON DISCAPACIDAD INTELECTUAL
% ==============================================================================
\fichaAlumno[%
    nombre={Alejandro ``Álex'' García},
    nombreCorto={ÁLEX G.},
    edad={19 años},
    icono={\faUserGraduate},
    foto={img/alex_garcia.jpg},
    badgeRol={Alumno PTVAL},
    badgeTarea={Tareas Rotativas},
    centro={CEE Purísima Concepción},
    cita={Necesito saber con imágenes qué tengo que hacer después y cuánto tiempo me queda para no ponerme nervioso ni perderme en los pasillos.},
    demographic={%
        \item \textbf{Edad y Género (Age / Gender):} 19 años, Varón.
        \item \textbf{Ubicación (Location):} Granada capital (domicilio familiar).
        \item \textbf{Perfil diagnóstico:} Discapacidad intelectual moderada (DI) con rasgos de TEA. Lenguaje funcional telegráfico. Motricidad física preservada.
    },
    computerExp={%
        \item \textbf{Habilidades informáticas:} Nivel básico táctil asistido. No usa ratón ni teclado.
        \item \textbf{Años de experiencia:} 5 años usando tablets en el aula.
        \item \textbf{Hábitos de uso:} Navegación por fotos de contactos, vídeos y juegos de emparejamiento.
        \item \textbf{Autonomía digital:} Buena en pantallas con pocas opciones claras.
    },
    productExp={%
        \item \textbf{Productos similares:} Paneles físicos de velcro (PECS) y app \textit{DictaPicto}.
        \item \textbf{Limitaciones actuales:} Las tarjetas físicas de papel se pierden; las apps habituales tienen mucho texto o menús complejos.
        \item \textbf{Tendencias de uso:} Gran adherencia a apoyos con imágenes y sonidos.
    },
    tasks={%
        \item \textbf{Tareas Primarias:} 
        1) Login tocando su foto $\rightarrow$ 2) Consultar la tarea rotativa asignada (reprografía, comedor, orden) $\rightarrow$ 3) Seguir los pasos guiados con imágenes $\rightarrow$ 4) Confirmar paso verde.
        \item \textbf{Tareas Secundarias:} Mirar cuánto falta para el recreo; pedir ayuda al tutor.
    },
    domainKnowledge={%
        \item \textbf{Comprensión del área:} Conoce las rutinas del colegio, las aulas y a los maestros.
        \item \textbf{Comprensión temporal:} Relativa visual (``ahora'', ``después''). No interpreta números de reloj convencional.
        \item \textbf{Destreza procedimental:} Ejecuta tareas escolares paso a paso de forma autónoma con apoyo de imágenes.
    },
    technology={%
        \item \textbf{Hardware:} Tablet Android.
        \item \textbf{Software:} Síntesis de voz.
        \item \textbf{Herramientas habituales:} Conexión Wi-Fi.
    },
    attitudes={%
        \item \textbf{Preferencias:} Refuerzos sonoros alegres (aplausos) y pantallas con una sola idea.
        \item \textbf{Miedo a la tecnología:} Nulo temor a tocar; ansiedad ante pitidos de error o bloqueos de la aplicación.
        \item \textbf{Valores:} Gran satisfacción por sentirse útil participando en las tareas del colegio.
    },
    learningStyle={%
        \item \textbf{Perfil de aprendizaje:} \textbf{Visual estricto} reforzado por canal \textbf{Auditivo (TTS)}.
        \item \textbf{Canales óptimos:} Pictogramas ARASAAC y fotos reales de profesores y aulas.
        \item \textbf{Metodología:} Modelado, ensayo guiado (\textit{learning by doing}) y refuerzo inmediato.
    },
    criticality={%
        \item \textbf{Consecuencias de un error:} \textit{Emocional:} Desorientación, ansiedad y bloqueo con llanto o rechazo. \textit{Operativo:} Confusión de aulas o retraso en las entregas.
        \item \textbf{Tolerancia del software:} Extrema. Sin acciones destructivas, confirmación asistida y botón de ayuda permanente.
    },
    escenario={A las 10:15h, la tablet emite una melodía suave. Álex toca su foto y descubre su tarea rotativa de la semana: reprografía. Sigue los pasos guiados con fotos de profesores, recoge los folios y pulsa el botón verde. Tras repartir las copias, la pantalla celebra su éxito con un aplauso sonoro y suma una estrella a su casillero.},
    badgesAjustes={%
        \badge{primaryColor}{\faKey\ Login Foto/Color} \hspace{2mm}
        \badge{secondaryColor}{\faIcons\ ARASAAC + Fotos Reales} \hspace{2mm}
        \badge{accentColor}{\faSync\ Tareas Rotativas} \hspace{2mm}
        \badge{purpleColor}{\faVolumeUp\ Síntesis de Voz TTS}%
    }%
]

% ==============================================================================
% PERSONA 3: MADRE DE UN NIÑO
% ==============================================================================
\fichaAdulto[%
    nombre={Carmen Fernández Delgado},
    nombreCorto={CARMEN F.},
    edad={46 años},
    icono={\faUserNurse},
    foto={img/carmen_fernandez.jpg},
    badgeRol={Madre de Alumno PTVAL},
    badgeTarea={Coordinación del Hogar},
    centro={Entorno Familiar \textbar\ Granada},
    cita={Quiero que la autonomía que Carlos gana en el colegio continúe en casa; necesito una agenda donde podamos programar sus rutinas de la tarde sin que suponga otra carga técnica inasumible tras mi jornada laboral.},
    demographic={%
        \item \textbf{Edad y Género:} 46 años, Mujer.
        \item \textbf{Ubicación:} Granada capital (barrio Zaidín).
        \item \textbf{Situación familiar:} Madre de Carlos (alumno con discapacidad motriz). Cuidadora principal.
    },
    occupation={%
        \item \textbf{Puesto actual:} Administrativa contable en empresa local (jornada continua).
        \item \textbf{Antigüedad:} 15 años en la empresa \textbar\ 18 años como madre y cuidadora.
        \item \textbf{Responsabilidades:} Conciliación laboral y familiar, coordinación de fisioterapia semanal y apoyo en el hogar.
    },
    education={%
        \item \textbf{Titulación (Degree):} Diplomatura en Ciencias Empresariales.
        \item \textbf{Especialidad (Major):} Gestión contable y financiera.
        \item \textbf{Cursos realizados:} Talleres de escuela de familias sobre productos de apoyo y autonomía en el hogar.
    },
    computerExp={%
        \item \textbf{Habilidades informáticas:} Nivel usuario medio en smartphone y ordenador.
        \item \textbf{Años de experiencia:} +16 años utilizando dispositivos digitales.
        \item \textbf{Hábitos de uso:} Uso diario de apps de mensajería, trámites online y banca móvil.
    },
    productExp={%
        \item \textbf{Productos similares:} Agendas de papel escolar, notas en la nevera y calendario del móvil.
        \item \textbf{Limitaciones actuales:} Los calendarios convencionales no están adaptados a su hijo y las notas en papel se pierden.
    },
    tasks={%
        \item \textbf{Tareas Primarias:} 
        1) Comprobar en el móvil las tareas escolares de Carlos $\rightarrow$ 2) Programar las rutinas de la tarde en casa $\rightarrow$ 3) Añadir fotos cotidianas a los pasos $\rightarrow$ 4) Supervisar el progreso.
        \item \textbf{Tareas Secundarias:} Enviar notas rápidas al tutor sobre estado de ánimo o incidencias.
    },
    domainKnowledge={%
        \item \textbf{Comprensión del área:} Conoce los tiempos de fatiga de su hijo y la importancia de fomentar su autonomía sin sobreprotegerlo.
        \item \textbf{Enfoque educativo:} Muy alineada con la transición a la vida adulta del centro escolar.
    },
    technology={%
        \item \textbf{Hardware:} Móvil Android y ordenador.
        \item \textbf{Software:} App móvil.
        \item \textbf{Herramientas:} Conexión a internet.
    },
    attitudes={%
        \item \textbf{Preferencias:} Interfaz rápida, directa, sin configuraciones farragosas y con sincronización en tiempo real.
        \item \textbf{Miedo a la tecnología:} Preocupación por caídas de conexión que desorienten a su hijo al llegar a casa.
        \item \textbf{Valores:} Bienestar emocional de su hijo e impulso a su independencia personal.
    },
    learningStyle={%
        \item \textbf{Perfil de aprendizaje:} Práctico y funcional; resolución directa de tareas cotidianas.
        \item \textbf{Canales óptimos:} Interfaz clara e intuitiva sin necesidad de manuales extensos.
        \item \textbf{Curva de adopción:} Inmediata si comprueba que su hijo usa la agenda en casa sin frustración.
    },
    criticality={%
        \item \textbf{Consecuencias de un error:} \textit{Logístico/Familiar:} Olvido de una rutina de fisioterapia o merienda. \textit{Emocional:} Estrés y sobrecarga tras la jornada laboral.
        \item \textbf{Tolerancia del software:} Guardado local ante fallos de cobertura y sincronización en la nube.
    },
    escenario={A las 16:30h al salir de la oficina, Carmen consulta la app en su móvil y comprueba que Carlos completó con éxito su tarea rotativa de la mañana. Con 3 toques rápidos programa la rutina de la tarde: ``17:30h - Merienda'' y ``18:15h - Preparar mochila de fisioterapia'' con fotos de la mochila azul. Carlos recibe las tareas en su tablet y Carmen se siente tranquila al ver cómo su hijo las completa.},
    badgesAjustes={%
        \badge{tealColor}{\faMobile*\ App Móvil Familiar} \hspace{2mm}
        \badge{primaryColor}{\faHome\ Rutinas del Hogar} \hspace{2mm}
        \badge{accentColor}{\faSync\ Sincronización Centro-Casa} \hspace{2mm}
        \badge{greenColor}{\faComments\ Chat Directo Tutor}%
    }%
]

% ==============================================================================
% PERSONA 4: PROFESOR (TUTOR ESCOLAR)
% ==============================================================================
\fichaAdulto[%
    nombre={Javier Morales Benítez},
    nombreCorto={JAVIER M.},
    edad={39 años},
    icono={\faChalkboardTeacher},
    foto={img/javier_morales.jpg},
    badgeRol={Tutor Docente PTVAL},
    badgeTarea={Coordinación de Aula},
    centro={CEE Purísima Concepción},
    cita={Gestiono un aula heterogénea con diferentes niveles de autonomía. Necesito organizar el día de cada alumno en minutos antes de las 9:00, asignar las tareas rotativas de reprografía o comedor, y tener la certeza de que el sistema avisa a cocina sin errores.},
    demographic={%
        \item \textbf{Edad y Género:} 39 años, Varón.
        \item \textbf{Ubicación:} Granada capital.
        \item \textbf{Perfil profesional:} Maestro especialista en Pedagogía Terapéutica (PT) con 12 años de experiencia docente, 5 en el CEE Purísima Concepción.
    },
    occupation={%
        \item \textbf{Puesto actual:} Tutor de Aula del Programa de Transición a la Vida Adulta (PTVAL).
        \item \textbf{Antigüedad:} 12 años en educación especial \textbar\ 5 años en el centro actual.
        \item \textbf{Responsabilidades:} Planificar itinerarios de autonomía, asignar roles de centro (reprografía, comedor), coordinar con familias y velar por la seguridad del aula.
    },
    education={%
        \item \textbf{Titulación (Degree):} Grado en Educación Primaria (Mención Pedagogía Terapéutica).
        \item \textbf{Especialidad (Major):} Pedagogía Terapéutica y Audición y Lenguaje (AL).
        \item \textbf{Cursos realizados:} Formación continua en TIC accesibles, sistemas SAAC y metodologías de inclusión laboral.
    },
    computerExp={%
        \item \textbf{Habilidades informáticas:} Nivel avanzado en herramientas educativas digitales y pictogramas ARASAAC.
        \item \textbf{Años de experiencia:} +15 años gestionando plataformas digitales y pizarras interactivas.
        \item \textbf{Hábitos de uso:} Uso continuo de ordenador portátil, tablet y smartphone.
    },
    productExp={%
        \item \textbf{Productos similares:} Portales de pictogramas, hojas de cálculo para el comedor y paneles físicos de aula.
        \item \textbf{Limitaciones actuales:} Pérdida de tiempo recortando velcros y falta de seguimiento en tiempo real fuera del aula.
    },
    tasks={%
        \item \textbf{Tareas Primarias:} 
        1) Configurar tareas diarias antes de la asamblea $\rightarrow$ 2) Asignar encargos rotativos del día a los alumnos $\rightarrow$ 3) Supervisar el estado del circuito $\rightarrow$ 4) Validar listado de comedor para cocina.
        \item \textbf{Tareas Secundarias:} Enviar mensajes de refuerzo y atender dudas por chat accesible.
    },
    domainKnowledge={%
        \item \textbf{Comprensión del área:} Experto en didáctica adaptada, diversidad cognitiva y motriz, sistemas SAAC y anticipación temporal.
        \item \textbf{Conocimiento del grupo:} Conoce el umbral de tolerancia a la frustración y ritmos de cada uno de sus alumnos.
    },
    technology={%
        \item \textbf{Hardware:} Ordenador portátil y tablet.
        \item \textbf{Software:} Navegador web.
        \item \textbf{Herramientas:} Red Wi-Fi.
    },
    attitudes={%
        \item \textbf{Preferencias:} Software ágil, paneles limpios tipo dashboard sin pasos redundantes y visión del aula en un vistazo.
        \item \textbf{Miedo a la tecnología:} Preocupación por caídas de red que dejen a los alumnos sin soporte visual durante la jornada.
        \item \textbf{Valores:} Compromiso con la inclusión y la capacitación laboral de sus alumnos.
    },
    learningStyle={%
        \item \textbf{Perfil de aprendizaje:} Analítico, estructurado y visual.
        \item \textbf{Canales óptimos:} Paneles de control con tarjetas de estado y tablas dinámicas.
        \item \textbf{Adopción:} Rápida si ahorra tiempo administrativo en la gestión del aula.
    },
    criticality={%
        \item \textbf{Consecuencias de un error:} \textit{Logístico y de Salud:} Notificar erróneamente dietas de comedor (alergias o celiaquía). \textit{Conductual:} Desajuste horario que provoque crisis de ansiedad en el alumno.
        \item \textbf{Tolerancia del software:} Confirmación obligatoria para dietas de comedor y registro histórico de cambios.
    },
    escenario={A las 8:40h antes de que lleguen los alumnos, Javier abre el panel web de CeRA desde su portátil. En dos minutos distribuye los encargos rotativos del día (reprografía a Álex, visitas a Carlos) con las fotos y apoyos asignados. A las 10:20h el tablero le muestra que las tareas avanzan y valida con un clic el listado de comensales para cocina.},
    badgesAjustes={%
        \badge{tealColor}{\faDesktop\ Panel Web Docente} \hspace{2mm}
        \badge{primaryColor}{\faTasks\ Asignación de Rutinas} \hspace{2mm}
        \badge{accentColor}{\faUtensils\ Control Dietas Comedor} \hspace{2mm}
        \badge{greenColor}{\faChartLine\ Monitorización Directo}%
    }%
]

\end{document}
